\section{The interstellar radiation field: photoelectric heating and
H$_2$ photodissociation}
\label{sec:fuv-isrf}

\subsection{Purpose}
\label{sec:fuv-purpose}

Massive stars emit most of their light below the hydrogen ionization
edge. That non-ionizing far-ultraviolet (FUV) radiation heats the gas by
ejecting electrons from dust grains (photoelectric heating) and
dissociates H$_2$ through the Lyman-Werner (LW) bands. Both processes
depend on the \emph{local} radiation field, which varies by orders of
magnitude between a star-forming region and the diffuse interstellar
medium. This module, part of the GEAR subgrid model
\citep{RevazJablonka2018} in \swift{} \citep{Schaller2024}, estimates that
field for every gas particle and hands it to the chemistry and cooling
solver \citep[Grackle,][]{Smith2017}.

\textbf{What the module provides.} For each gas particle $i$ it carries
two radiation energies per unit gas mass, $u_{{\rm FUV},i}$ for the
$6$--$11.2$\,eV band (the PE band of \citealt{Kim2023}) and $u_{{\rm LW},i}$ for the $11.2$--$13.6$\,eV band, in
${\rm erg\,g^{-1}}$, sourced by the stars in the simulation and
attenuated by dust. The code stores the two bands under the internal
identifiers \code{ISRF\_BAND\_PE} and \code{ISRF\_BAND\_LW}; this document
keeps calling the first band FUV throughout, and the snapshot I/O fields
keep the FUV name too (\code{FUVSpecificEnergy}, \code{LWSpecificEnergy}).
From them it derives, at every cooling call, the
photoelectric heating strength $G_0$ in Habing units and the H$_2$
photodissociation rate $k_{\rm diss}$, which Grackle consumes through two
of its own per-cell input fields (Section~\ref{sec:fuv-coupling}).

\textbf{Volumetric or per unit mass?} The transport is written in
Lagrangian form on the gas particles, so the two evolved variables of each
band are both per unit gas mass: the radiation energy $u_b$
(${\rm erg\,g^{-1}}$) and the flux $\mathbf F_b = \mathbf F_{V,b}/\rho$,
with $\rho$ the gas density. Those are the quantities the particles store
and the quantities every update below advances. The volumetric energy
density is recovered from them as $U_b = \rho\,u_b$
(${\rm erg\,cm^{-3}}$), and it is the volumetric field, through the energy
flux $c\,U_b$, that the chemistry reads to form $G_0$ and $k_{\rm diss}$.
Both forms appear below, and which one is meant is always visible in the
symbol: $u_b$ and $\mathbf F_b$ wherever the scheme evolves, stores or
exchanges a field, $U_b$ and $\mathbf F_{V,b}$ wherever a continuum moment
equation, an analytic profile or the chemistry coupling is written. The
word \emph{specific} is used here only in the sense of per unit mass.

\textbf{Two run modes.} With propagation off (the default), each star
deposits its band luminosities on the gas inside its own smoothing
kernel and nothing else happens: the field is an instantaneous snapshot
of the illumination within one kernel radius of each star
(Section~\ref{sec:fuv-instantaneous}). With propagation on, the deposit
is the source term of a two-moment (energy and flux) radiation transport
system, solved on the SPH particles with a reduced signal speed, so that
the field reaches beyond the kernel, is smooth between sources, and decays
over the dust absorption length (Sections~\ref{sec:fuv-equations} and
\ref{sec:fuv-numerics}).

\textbf{What the module deliberately does not do.} It is not a full
radiative transfer solver: two angular moments only, no scattering, no
frequency resolution inside a band, and a signal speed far below $c$.
It carries no ionizing band (photoionization is the separate HII model)
and exerts no momentum on the gas (radiation pressure is the separate
channel of Section~\ref{sec:radiation-pressure}). It produces only the
stellar-sourced local field. For H$_2$ dissociation this local rate is
added to Grackle's own uniform, time-dependent metagalactic background
rate and the sum is self-shielded together
(Section~\ref{sec:fuv-coupling}). For photoelectric heating and the
other dust channels there is no such addition: the per-particle field
replaces Grackle's scalar background outright, so an unilluminated
particle carries no residual background heating. H$_2$ self-shielding is
not modelled here: Grackle attenuates the supplied dissociation rate
with its own self-shielding prescription \citep{WolcottGreenHaiman2019},
exactly as it attenuates the background rate.

\subsection{The equations we solve}
\label{sec:fuv-equations}

\subsubsection{Fields, bands and units}
\label{sec:fuv-fields}

Each band $b\in\{{\rm FUV},{\rm LW}\}$ is described by its volumetric
energy density $U_b$ (${\rm erg\,cm^{-3}}$) and its flux $\mathbf
F_{V,b}$ (${\rm erg\,cm^{-2}\,s^{-1}}$), the first two angular moments of
the specific intensity. The particles carry the mass-specific
counterparts,
\begin{equation}
  u_b = \frac{U_b}{\rho}\ \ ({\rm erg\,g^{-1}}), \qquad
  \mathbf F_b = \frac{\mathbf F_{V,b}}{\rho}\ \ ({\rm erg\,cm\,g^{-1}\,s^{-1}}),
  \label{eq:fuv-specific-fields}
\end{equation}
with $\rho$ the gas density, in the same way SPH carries a specific
rather than a volumetric internal energy. The band edges are the
classical Habing range \citep{Habing1968}, split at the lower edge of the
Lyman-Werner bands: photoelectric heating responds to the whole $6$--$13.6$\,eV range,
H$_2$ photodissociation only to the harder $11.2$--$13.6$\,eV part. Both
bands obey the same equations with band-specific opacities and sources;
the band index is dropped below whenever no confusion is possible.

\subsubsection{Sources: stellar band luminosities}
\label{sec:fuv-sources}

A star particle emits $L_{\rm FUV}$ and $L_{\rm LW}$ (${\rm erg\,s^{-1}}$),
the integrals of its spectrum over the two bands of
Section~\ref{sec:fuv-fields}, $6$--$11.2$\,eV and $11.2$--$13.6$\,eV. Both
are tabulated per unit initial stellar mass, in the same stellar tables that
supply the bolometric luminosity and the ionizing photon rate, and are read
by interpolation on those tables' own axes.

An individual star reads an entry indexed by its initial mass and its
metallicity. The value is a main-sequence average and not a function of
age, the same convention $L_{\rm bol}$ follows. A stellar-population
particle reads an IMF-integrated band luminosity for its metallicity and
its current age, the age entering through the surviving mass range of the
integral, and multiplies it by the particle's initial mass; where the table
carries a main-sequence lifetime, that range is capped at the mass leaving
the main sequence at the current age, so stars past the main sequence
contribute to neither band and the integrated band luminosities decline
with age. A table that carries no band data at all is covered by a
fallback, which splits the two bands off the bolometric luminosity using
the star's tabulated effective temperature.

Both luminosities are deposited on the gas particles inside the star's own
kernel, kernel-weighted and attenuated by dust
(Section~\ref{sec:fuv-injection}). A star that is no longer an active
feedback source makes no further deposits in either band; that, rather than
any change to the tabulated values, is how a source switches off
(Section~\ref{sec:fuv-turnoff}). The tabulated band luminosities are
deposited as they stand, with no Doppler correction for the star's motion
relative to the gas it illuminates: that correction shifts each band edge
by the fraction $v/c$, below $3\times10^{-4}$ for any peculiar velocity in
these simulations, the same order as the band-edge term the transport drops
(Section~\ref{sec:fuv-assumptions}).

\subsubsection{Dust absorption}
\label{sec:fuv-opacity}

Dust is the only absorber. The band-specific dust mass opacity is
\begin{equation}
  \kappa_b(Z) = \frac{\sigma_{d,b}}{\mu_H m_H}\,\mathcal D(Z), \qquad
  \mathcal D(Z) = \frac{Z}{Z_\odot}\,\frac{f_{\rm gr}}{f_{\rm gr,MW}},
  \label{eq:fuv-kappa}
\end{equation}
with $\sigma_{d,{\rm FUV}} = 9\times10^{-22}\,{\rm cm^2}$ and
$\sigma_{d,{\rm LW}} = 1.5\times10^{-21}\,{\rm cm^2}$ the dust
absorption cross-sections per hydrogen nucleon in the two bands,
band-averaged over a coeval, $\sim$50\,Myr, Kroupa-IMF stellar
population for Milky-Way grains \citep{WeingartnerDraine2001},
matching \citet[their PE and LW bands, Section~2.4]{Kim2023} exactly,
$\mu_H = 1.4$ the mean mass per hydrogen nucleon in units of
$m_H$, and $\mathcal D(Z)$ the dust-to-gas ratio relative to the Milky
Way. $\mathcal D$ scales linearly with the gas metal mass fraction $Z$,
normalised to Grackle's solar value $Z_\odot = 0.01295$ and to Grackle's
default dust-to-gas ratio $f_{\rm gr,MW} = 0.009387$, so that this module
and Grackle's dust chemistry assume the same dust abundance for the same
gas. At solar metallicity $\kappa_{\rm FUV}\approx 380$ and $\kappa_{\rm
LW}\approx 640\,{\rm cm^2\,g^{-1}}$. The absorption length and the
absorption time at the true speed of light are
\begin{equation}
  \lambda_b = \frac{1}{\kappa_b\rho}, \qquad
  \tau_{b,{\rm phys}} = \frac{\lambda_b}{c} .
  \label{eq:fuv-lambda}
\end{equation}
$h/\lambda$, the ratio of the smoothing length to the absorption length,
is the resolution parameter that controls the behaviour of every part of
the scheme; it ranges from $\sim10^{-3}$ in diffuse, metal-poor gas to
$\sim20$ in dense, metal-rich gas at production resolution.

\subsubsection{Moment equations and closure}
\label{sec:fuv-moments}

The first two angular moments of the radiative transfer equation in a
purely absorbing medium with an emissivity $S_V$ (${\rm
erg\,cm^{-3}\,s^{-1}}$) are
\begin{align}
  \frac{\partial U}{\partial t} + \nabla\cdot\mathbf F_V
    &= -c\,\kappa\rho\,U + S_V ,
  \label{eq:fuv-moment-zero} \\
  \frac{\partial \mathbf F_V}{\partial t} + c^2\,\nabla\cdot\mathbf P
    &= -c\,\kappa\rho\,\mathbf F_V ,
  \label{eq:fuv-moment-one}
\end{align}
where $\mathbf P$ is the radiation pressure tensor. The system is closed
with the M1 closure \citep{Levermore1984,DubrocaFeugeas1999},
\begin{equation}
  \mathbf P = \mathbf D(f)\,U , \qquad
  \mathbf D(f) = \frac{1-\chi}{2}\,\mathbf I + \frac{3\chi-1}{2}\,
    \hat{\mathbf n}\otimes\hat{\mathbf n} , \qquad
  \chi(f) = \frac{3+4f^2}{5+2\sqrt{4-3f^2}} ,
  \label{eq:fuv-m1-closure}
\end{equation}
with $\hat{\mathbf n} = \mathbf F/|\mathbf F|$ the flux direction and $f
= |\mathbf F_V|/(c\,U)\in[0,1]$ the reduced flux. At $f = 0$, $\mathbf D
= \mathbf I/3$ (isotropic radiation, the Eddington or P1 closure); at $f
= 1$, $\mathbf D = \hat{\mathbf n}\otimes\hat{\mathbf n}$ (free streaming
along $\hat{\mathbf n}$). The reduced flux is the one quantity of the
closure that does not care which of the two variable sets it is built from:
$|\mathbf F_V|/(c\,U) = |\mathbf F|/(c\,u)$ identically, since the density
cancels. The moment equations do not by themselves
forbid $|\mathbf F_V| > cU$; the admissibility constraint
\begin{equation}
  |\mathbf F_V| \le c\,U
  \label{eq:fuv-admissible}
\end{equation}
is imposed explicitly (Section~\ref{sec:fuv-limiter}).

\textbf{Precision of the reduced flux.} The squared flux norm that $f$
above, the flux limiter of Eq.~\ref{eq:fuv-flux-limiter} and the
flux-relaxation-residual gate of Eq.~\ref{eq:fuv-diss-gate} all read
from $\mathbf F$ is evaluated in double precision, not the particle's
native float32. A well-attenuated far band routinely carries
$|\mathbf F|$ below $\sqrt{\code{FLT\_MIN}}\approx1.08\times10^{-19}$ in
the code's internal units, and squaring a value that small in float32
underflows silently to exactly zero: the closure then reads $f=0$ and
returns the isotropic $\mathbf D=\mathbf I/3$ even though the physical
flux is merely small rather than zero, the limiter's clamp is skipped
because it sees no flux to limit, and the gate's residual denominator
vanishes with it. Only the squared norm, its reciprocal square root and
the gate's residual ratio are widened; the closure coefficients,
$\mathbf D$ itself and $u$ stay float32.

\textbf{Reduced signal speed.} The characteristic speeds of
Eqs.~\ref{eq:fuv-moment-zero}--\ref{eq:fuv-moment-one} lie in
$[-c, c]$, which would impose a time step $\Delta t \lesssim h/c$, far
below any hydrodynamical or cooling time step. The equations are
therefore solved with the speed of light replaced by a reduced speed
$c_{\rm hyp}\ll c$ everywhere it appears (absorption rate, pressure
term, reduced flux, admissibility constraint), and with the source
rescaled by $c_{\rm hyp}/c$ so that the steady-state field is unchanged:
the reduced-speed-of-light approximation
\citep{GnedinAbel2001,Rosdahl2013}. The moments are carried on the SPH
particles themselves rather than on a grid, in the manner of the SPH M1
scheme of \citet{Chan2021}.

Eqs.~\ref{eq:fuv-moment-zero}--\ref{eq:fuv-moment-one} are already a
reduced form, and the next section states what the reduction assumes.
Two further changes then take them into the form the code integrates: the
variables per unit gas mass of Eq.~\ref{eq:fuv-specific-fields}, and the
expanding, gas-following frame (Section~\ref{sec:fuv-lagrangian}).

\subsubsection{Assumptions and neglected terms}
\label{sec:fuv-assumptions}

Written in the frame of the gas, which is the frame in which the dust
opacity of Section~\ref{sec:fuv-opacity} is the tabulated one, and kept to
first order in $v/c$ for one grey band, the complete pair of moment
equations is
\begin{align}
  \frac{\partial U}{\partial t}
    &+ \underbrace{\nabla\cdot(U\mathbf v)}_{\text{advection}}
    + \underbrace{\nabla\cdot\mathbf F_V}_{\text{transport}}
    + \underbrace{\mathbf P\!:\!\nabla\mathbf v}_{\text{compression work}}
  \nonumber \\
  &= \underbrace{S_V}_{\text{sources}}
    + \underbrace{4\pi\kappa_P\rho\,B}_{\text{thermal emission}}
    - \underbrace{c\,\kappa_a\rho\,U}_{\text{absorption}}
    + \underbrace{\Delta_{\rm edge}}_{\text{band-edge shift}}
    - \underbrace{3HU}_{\text{volume}}
    - \underbrace{HU}_{\text{redshift}} ,
  \label{eq:fuv-full-moment-zero} \\
  \frac{\partial \mathbf F_V}{\partial t}
    &+ \underbrace{\nabla\cdot(\mathbf F_V\otimes\mathbf v)}_{\text{advection}}
    + \underbrace{(\mathbf F_V\cdot\nabla)\mathbf v}_{\text{aberration}}
    + \underbrace{c^2\,\nabla\cdot\mathbf P}_{\text{transport}}
  \nonumber \\
  &= -\underbrace{c\,(\kappa_a+\kappa_s)\rho\,\mathbf F_V}
        _{\text{absorption, scattering}}
    + \underbrace{\Delta_{\rm edge,F}}_{\text{band-edge shift}}
    - \underbrace{3H\mathbf F_V}_{\text{volume}}
    - \underbrace{H\mathbf F_V}_{\text{redshift}} .
  \label{eq:fuv-full-moment-one}
\end{align}
Here $\kappa_a$ and $\kappa_s$ are the absorption and scattering parts of
the dust opacity, $B$ the Planck function at the dust temperature with
$\kappa_P$ its Planck mean, $\mathbf v$ the gas velocity, and $\Delta_{\rm
edge}$ the net transfer of energy across the two edges of a band of finite
width, which replaces the frequency-derivative term of the monochromatic
equations: a band edge sits at a different frequency in the gas frame than
in the frame of the source, by the fractional amount $v/c$, so energy
crosses it. The exact coefficients of the $v/c$ terms in the first moment
follow the standard comoving-frame reduction
\citep{MihalasMihalas1984}, in the mixed-frame ordering of
\citet{Lowrie1999}; the moment schemes this one follows make the same
neglect \citep{Rosdahl2013,Chan2021}. Coherent isotropic scattering does not appear in the
energy equation at leading order, because what it removes from a direction
it returns from another; it removes flux at the total opacity, which is why
$\kappa_a+\kappa_s$ appears in the first moment and $\kappa_a$ alone in the
zeroth.

The ordering that follows is done on the physical system, at the true speed
of light. Every estimate compares a term of
Eqs.~\ref{eq:fuv-full-moment-zero}--\ref{eq:fuv-full-moment-one} against a
retained one, and the smallness of the dropped terms is a property of the
photons themselves: of their real travel, over the real absorption length
$\lambda$, in the real time $\lambda/c$. The reduced signal speed is
introduced afterwards, in the transport and the source together, and leaves
the steady state of what remains unchanged (Section~\ref{sec:fuv-steady}).
Estimating the dropped terms against $c_{\rm hyp}$ instead would inflate
them by $c/c_{\rm hyp}$ without any photon travelling differently.

\textbf{What we drop, and why.}
\begin{enumerate}
  \item \textbf{Aberration of the flux}, $(\mathbf F_V\cdot\nabla)\mathbf
    v$. The term tilts the flux direction by an angle of order $v/c$.
    Peculiar velocities of $10$ to $100\,{\rm km\,s^{-1}}$ give $v/c =
    3\times10^{-5}$ to $3\times10^{-4}$, so the tilt is at most a
    thousandth of a radian, orders of magnitude below the directional
    error a two-moment closure already makes when it merges two crossing
    beams into one direction.
  \item \textbf{The band-edge Doppler shift}, $\Delta_{\rm edge}$ and
    $\Delta_{\rm edge,F}$. Each edge moves by the fractional amount
    $v/c$, at most $3\times10^{-4}$ for the velocities above, and the
    band-integrated energy changes by that same fraction: four orders of
    magnitude below the accuracy of the dust cross-sections the band
    itself is built on. This is why the transport carries no Doppler
    correction anywhere, on the fields or on the sources
    (Section~\ref{sec:fuv-sources}). Where the Doppler shift does matter
    is one level down, inside the LW band: the H$_2$ absorption lines are
    narrow, their widths set by the Doppler parameter $b$ of a few
    ${\rm km\,s^{-1}}$ from thermal and turbulent motions, which is
    comparable with the relative velocities themselves, so a shift of
    that size moves a line on and off the part of the spectrum that
    shields it. That physics is not transport and cannot be represented
    by a band-integrated energy: it enters through the $b$ dependence of
    the H$_2$ self-shielding function that the chemistry applies
    \citep{WolcottGreenHaiman2019}, to the rate this module supplies
    (Section~\ref{sec:fuv-coupling}). Splitting the band into lines is
    the only way to do better, and it is outside a two-moment grey
    scheme.
  \item \textbf{The compression work term}, $\mathbf P\!:\!\nabla\mathbf
    v$. Against the absorption term, its ratio is $|\mathbf
    P\!:\!\nabla\mathbf v|/(c\kappa_a\rho\,U)\lesssim
    |\nabla\mathbf v|\,\lambda/c$, the free-streaming contraction being at
    most three times the isotropic value $U/3$. A velocity of
    $10\,{\rm km\,s^{-1}}$ over $100$\,pc gives $|\nabla\mathbf v|\sim
    3\times10^{-15}\,{\rm s^{-1}}$, and $\lambda$ runs from $\sim0.1$\,pc
    in dense metal-rich gas to tens of kiloparsecs in diffuse metal-poor
    gas, so the ratio is $10^{-8}$ at one end and at most $10^{-2}$ at the
    other. The energy densities alone do not settle it: at $G_0 = 1$ the
    band energy density is $5\times10^{-14}\,{\rm erg\,cm^{-3}}$, three
    per cent of the thermal energy density of warm neutral gas at
    $n_H = 1\,{\rm cm^{-3}}$ and $T = 8000$\,K, but comparable with that
    of cold neutral gas at the same density and $T = 100$\,K, and larger
    again inside a bright photodissociation region. What closes the item
    is that this module applies no force and no work to the gas at all:
    momentum coupling is the separate channel of
    Section~\ref{sec:radiation-pressure}. The radiation therefore does no
    $P\,{\rm d}V$ work on the gas and receives none from compression, and
    the dropped term could only have acted on the radiation's own budget,
    where the bound above holds.
  \item \textbf{Dust scattering}, $\kappa_s$. The tabulated values of
    Section~\ref{sec:fuv-opacity} are absorption cross-sections, so
    $\kappa_a = \kappa$ and $\kappa_s = 0$, and the single opacity of
    Eqs.~\ref{eq:fuv-moment-zero}--\ref{eq:fuv-moment-one} is the
    consistent choice. Were extinction values used instead, with a
    Milky-Way grain albedo of a few tenths in the FUV, the flux would have
    to relax at the total opacity while the energy equation kept only the
    absorbing part; using extinction as if it were absorption instead, as
    a single opacity in both, destroys the energy that scattering returns
    to the diffuse field, and the absorption length used would be shorter
    than the true one by the factor $1-\text{albedo}$, between about
    $1.4$ and $2$ times too short.
    A missing scattering source also makes shadows behind dense clumps
    emptier than reality, in the opposite direction to the softening a
    two-moment closure produces.
  \item \textbf{Spectral migration across the band edges from expansion.}
    A photon lives, before absorption, for $\tau_{\rm phys} = \lambda/c$,
    which is $10^3$ to $10^5$\,yr over the $\lambda \sim 300$\,pc to
    $30$\,kpc of diffuse gas, against a Hubble time of order
    $10^{10}$\,yr. Its frequency therefore migrates by
    $H\tau_{\rm phys}\sim 10^{-7}$ to $10^{-5}$ of a band width before it
    is absorbed, and no band-edge bookkeeping is needed. The energy
    redshift itself is kept, as the $-Hu$ term of
    Section~\ref{sec:fuv-lagrangian}, whose Step~2 records the sense in
    which that term bounds the loss from a narrow band.
  \item \textbf{The true speed of light in the transport.} Replacing $c$
    by $c_{\rm hyp}$ and rescaling the source by $c_{\rm hyp}/c$ leaves
    every steady state of the system unchanged
    (Section~\ref{sec:fuv-steady}) and changes only how long the transient
    takes to settle: the standard argument for the approximation
    \citep{GnedinAbel2001,Rosdahl2013}, valid as long as the light
    crossing time of the region of interest stays short against the time
    on which the sources and the absorbers change. It fails where a front
    would move at a speed approaching $c_{\rm hyp}$, and for a uniform
    background rather than local sources \citep{Gnedin2016}. Here the
    cost is a transient longer than the lifetime of the star that drives
    it in the most diffuse gas, quantified in
    Section~\ref{sec:fuv-chyp}.
  \item \textbf{Thermal dust emission}, $4\pi\kappa_P\rho\,B$. Dust
    re-emits the FUV energy it absorbs in the infrared, far below the
    $6$\,eV lower edge, so no emission term returns energy to either
    band and the emissivity retains the stellar sources only. Both bands
    are optically thin to their own re-emission because there is none.
\end{enumerate}

\textbf{What we keep.} Absorption by dust at the tabulated opacity, in
both equations. The stellar sources of Section~\ref{sec:fuv-sources}. The
advection of both fields by the gas, exactly and with no discrete
operator, which is what moving the variables onto the gas particles buys
(Section~\ref{sec:fuv-lagrangian}). Both expansion terms, the volume
dilution and the redshift, one power of $H$ surviving in each equation
once the fields are measured per unit gas mass. And the M1 closure with
its admissibility constraint, which is what carries the directionality
that two moments can carry.

\subsubsection{From the Eulerian moments to the system we integrate}
\label{sec:fuv-lagrangian}

Eqs.~\ref{eq:fuv-moment-zero}--\ref{eq:fuv-moment-one} are volumetric,
physical and Eulerian: they hold at a fixed point of a static space. The
scheme integrates them per unit gas mass, on particles that move with the
gas, in a box that may expand. The mechanism comes first, in a static box;
five steps then restore the expansion terms.

\textbf{What a particle carries, and which derivative advances it.} Each
gas particle carries the two mass-specific fields of one band, $u$ and
$\mathbf F$, and it moves with the gas. The rate of change it experiences
is therefore not the rate at a fixed point of space but the rate along its
own trajectory, the material derivative
\begin{equation}
  \frac{{\rm D}}{{\rm D}t} = \frac{\partial}{\partial t} +
    \mathbf v\cdot\nabla ,
  \label{eq:fuv-lagrangian-derivative}
\end{equation}
with $\mathbf v$ the gas velocity. Everything the rest of this paragraph
does is turn the volumetric, fixed-point left-hand side of
Eq.~\ref{eq:fuv-full-moment-zero} into that derivative. Substitute $U =
\rho u$ and $\mathbf F_V = \rho\mathbf F$ into it, keeping absorption and
the source and taking $H = 0$:
\begin{equation}
  \frac{\partial(\rho u)}{\partial t}
  + \nabla\cdot\big(\rho\,u\,\mathbf v\big)
  + \nabla\cdot\big(\rho\,\mathbf F\big)
  = \rho\,\dot s - c\,\kappa\rho\,(\rho u) .
  \label{eq:fuv-lagr-line1}
\end{equation}
The gas obeys mass conservation, $\partial\rho/\partial t +
\nabla\cdot(\rho\mathbf v) = 0$, and for any quantity $\psi$ carried per
unit gas mass that identity reads $\partial(\rho\psi)/\partial t +
\nabla\cdot(\rho\psi\mathbf v) = \rho\,{\rm D}\psi/{\rm D}t$: the first
two terms of Eq.~\ref{eq:fuv-lagr-line1} are exactly $\rho\,{\rm D}u/{\rm
D}t$. Dividing by $\rho$,
\begin{equation}
  \frac{{\rm D}u}{{\rm D}t}
  = -\frac{1}{\rho}\nabla\cdot\big(\rho\,\mathbf F\big)
    - c\,\kappa\rho\,u + \dot s .
  \label{eq:fuv-lagr-line2}
\end{equation}
The advection term has not been neglected: it is the term that turned
$\partial/\partial t$ into ${\rm D}/{\rm D}t$, and the particle's own
motion supplies it, at no cost and with no discrete operator. Starting
instead from the fixed-point form
Eq.~\ref{eq:fuv-moment-zero}, which carries no advection term, the same
substitution would leave the residual
$+(1/\rho)\nabla\cdot(U\mathbf v)$ on the right of
Eq.~\ref{eq:fuv-lagr-line2}: the two forms differ by exactly that term, of
order $v/c$ relative to the transport term, and moving with the gas is
what makes the gas-frame form the one being solved rather than an
approximation to it. The other $v/c$ terms stay dropped, item by item in
Section~\ref{sec:fuv-assumptions}.

The first moment goes the same way, on each component of the flux:
\begin{equation}
  \frac{\partial(\rho \mathbf F)}{\partial t}
  + \nabla\cdot\big(\rho\,\mathbf F\otimes\mathbf v\big)
  + c^2\,\nabla\cdot\big(\mathbf D(f)\,\rho u\big)
  = -c\,\kappa\rho\,(\rho\mathbf F)
  \quad\Longrightarrow\quad
  \frac{{\rm D}\mathbf F}{{\rm D}t}
  = -\frac{c^2}{\rho}\nabla\cdot\big(\mathbf D(f)\,\rho u\big)
    - c\,\kappa\rho\,\mathbf F .
  \label{eq:fuv-lagr-F}
\end{equation}
Two spatial operators survive, and they are the only two the scheme has to
estimate: $(1/\rho)\nabla\cdot(\rho\mathbf F)$, the divergence of the
volumetric flux per unit density, and
$(c^2/\rho)\nabla\cdot(\mathbf D\,\rho u)$, the divergence of the
volumetric pressure tensor per unit density. Note what sits inside each
divergence: the density, because what crosses a surface is the volumetric
flux $\rho\mathbf F$ and the volumetric pressure $\mathbf D\,\rho u$, not
their mass-specific counterparts. Each is estimated by a sum over kernel
neighbours built so that the pairwise contribution to one particle is the
negative of the contribution to the other, weighted by their masses; the
transport then moves energy between particles and creates none, conserving
$\sum_i m_i u_i$ (Section~\ref{sec:fuv-operators}).

\textbf{Conventions.} \swift{} stores particle positions as comoving
coordinates $\mathbf x$, the physical position being $a\mathbf x$ with $a$
the scale factor, and gas densities as comoving densities $\rho_c$. Time is
proper time. The radiation variables are physical throughout. Hence
\begin{equation}
  \rho = a^{-3}\rho_c , \qquad
  \nabla = a^{-1}\nabla_{\mathbf x} , \qquad
  \frac{{\rm d}\ln a}{{\rm d}t} = H .
  \label{eq:fuv-cosmo-conventions}
\end{equation}

\textbf{Step 1: what expansion does to the volumetric field.} Written at a
fixed comoving position, the zeroth moment
Eq.~\ref{eq:fuv-moment-zero} acquires two expansion terms,
\begin{equation}
  \left.\frac{\partial U}{\partial t}\right|_{\mathbf x}
  + \underbrace{3HU}_{\text{volume}} + \underbrace{HU}_{\text{redshift}}
  + \frac{1}{a}\nabla_{\mathbf x}\cdot\mathbf F_V
  = -c\,\kappa\rho\,U + S_V ,
  \label{eq:fuv-cosmo-U}
\end{equation}
of different origin: $3HU$ is the dilution of a fixed energy over a
comoving volume growing as $a^3$, and $HU$ is the redshift of each photon
in it, $E\propto a^{-1}$. Together they give the textbook free-streaming
result, $U\propto a^{-4}$ with no source, no absorption and no net flux
divergence.

\textbf{Step 2: divide by the gas density.} The gas dilutes too, three
powers of it, so at fixed comoving gas density
\begin{equation}
  \left.\frac{\partial u}{\partial t}\right|_{\mathbf x}
  = \frac{1}{\rho}\left.\frac{\partial U}{\partial t}\right|_{\mathbf x}
    + 3H\,u ,
  \label{eq:fuv-cosmo-divide}
\end{equation}
and the three volume powers of Step 1 cancel against it exactly. One power
survives, the redshift one:
\begin{equation}
  \left.\frac{\partial u}{\partial t}\right|_{\mathbf x}
  + \frac{1}{a\,\rho}\nabla_{\mathbf x}\cdot\mathbf F_V
  = -H\,u - c\,\kappa\rho\,u + \dot s .
  \label{eq:fuv-cosmo-u}
\end{equation}
This is why the energy equation carries $-Hu$ and not $-4Hu$: dividing by
a gas density that dilutes as $a^{-3}$ has already absorbed the volume
part, and the field per unit gas mass loses only the redshift. Photons
redshifting out of a band across its edges are a further loss that $-Hu$
does not describe, so the term is a lower bound on the true loss from a
narrow band, not an overestimate. In a non-cosmological run $H=0$ and every
expansion term above vanishes.

\textbf{Step 3: the flux equation picks up its own $H$.} The first moment
carries the same four powers, most directly seen in the free-streaming
limit where $\mathbf F_V = c\,U\,\hat{\mathbf n}$ and therefore
$|\mathbf F_V|\propto a^{-4}$ as well; dividing by $\rho$ again removes
three and leaves one,
\begin{equation}
  \left.\frac{\partial \mathbf F}{\partial t}\right|_{\mathbf x}
  + \frac{c^2}{a\,\rho}\nabla_{\mathbf x}\cdot\big(\mathbf D\,U\big)
  = -H\,\mathbf F - c\,\kappa\rho\,\mathbf F .
  \label{eq:fuv-cosmo-F}
\end{equation}
Both equations therefore relax at the same rate, absorption plus
expansion, which is what lets one relaxation depth serve both.

\textbf{Step 4: follow the gas.} This is the substitution of
Eqs.~\ref{eq:fuv-lagr-line1}--\ref{eq:fuv-lagr-line2} above, applied now to
the two equations that carry the expansion terms: replace
$\partial/\partial t|_{\mathbf x}$ by ${\rm D}/{\rm D}t$ in
Eqs.~\ref{eq:fuv-cosmo-u} and \ref{eq:fuv-cosmo-F}, the advection of both
fields by the gas being accounted for by the particles' own motion. The
expansion terms are untouched by it: $H$ multiplies the local field, not a
gradient of it.

\textbf{Step 5: the transport operators the mass-specific form implies.}
These are the two operators already identified above, written once more
with the volumetric fields visible inside each divergence,
\begin{equation}
  \frac{1}{\rho}\nabla\cdot\mathbf F_V
    = \frac{1}{\rho}\nabla\cdot\big(\rho\,\mathbf F\big) , \qquad
  \frac{c^2}{\rho}\nabla\cdot\big(\mathbf D\,U\big)
    = \frac{c^2}{\rho}\nabla\cdot\big(\mathbf D\,\rho u\big) ,
  \label{eq:fuv-mass-specific-operators}
\end{equation}
and it is these two forms, rather than a bare divergence of $\mathbf F$ and
a bare gradient of $u$, that Section~\ref{sec:fuv-operators} discretises.

\textbf{The system integrated.} Collecting Steps 1 to 5 and rescaling by
$c_{\rm hyp}/c$: the transport and pressure-tensor divergences already
acquire that factor (once, respectively twice) from the plain
$c\to c_{\rm hyp}$ replacement inside them, but the absorption-plus-expansion
decay rate and the source term must be rescaled explicitly, as the whole
true-speed rate $c/\lambda + H = c\,\kappa\rho + H$, because $H$ carries no
factor of $c$ for the substitution to act on -- it is a property of the
expanding background, not of the radiation transport, so unlike
$\kappa$ it does not pick up the dilation "for free". Leaving $H$ out of
that rescaling, as an earlier version of this document did, is the same
class of error as leaving any other rate undilated:
\begin{align}
  \frac{{\rm D}u}{{\rm D}t}
    &= -\frac{1}{\rho}\nabla\cdot(\rho\,\mathbf F)
       - \frac{c_{\rm hyp}}{c}\left(\frac{c}{\lambda} + H\right) u
       + \frac{c_{\rm hyp}}{c}\,\dot s ,
  \label{eq:fuv-u-eq} \\
  \frac{{\rm D}\mathbf F}{{\rm D}t}
    &= -\frac{c_{\rm hyp}^2}{\rho}\nabla\cdot\big(\mathbf D(f)\,\rho u\big)
       - \frac{c_{\rm hyp}}{c}\left(\frac{c}{\lambda} + H\right)\mathbf F ,
  \qquad
  f = \min\!\left(\frac{|\mathbf F|}{c_{\rm hyp}u},\,1\right) ,
  \label{eq:fuv-F-eq}
\end{align}
with $\dot s = S_V/\rho$ the source rate per unit gas mass (${\rm
erg\,g^{-1}\,s^{-1}}$), $H$ the Hubble rate, and $c/\lambda = c\,\kappa\rho$
the true-speed absorption rate. This is the exact system the code
integrates. The gradients in it are physical; the discrete operators are
evaluated on comoving inputs and converted once, as
Section~\ref{sec:fuv-cosmo} shows. Both equations carry the same single
power of $H$, dilated by the same $c_{\rm hyp}/c$ factor as absorption
and injection, inside the same relaxation rate
$(c_{\rm hyp}/c)\big(c/\lambda + H\big)$, so absorption and expansion add
as rates rather than as timescales, and $c_{\rm hyp}$ cancels out of the
homogeneous equation's fixed point identically to every other term, for
any per-particle speed field (Section~\ref{sec:fuv-time}).

\textbf{Consequence of the undilated form.} Writing $H$ bare instead,
$(c_{\rm hyp}/\lambda + H)$, as this document did until this fix and as
the code did until \code{d47419d04}, suppresses the true-speed fixed
point $u^\ast_{\rm true} = (\dot s/c)/(\kappa\rho + H/c)$ by the factor
$x/(1+x)$, $x = c_{\rm hyp}\kappa\rho/H$: since $c_{\rm hyp}\ll c$,
$H/c_{\rm hyp} \gg H/c$, so the bare-$H$ fixed point
$(\dot s/c)/(\kappa\rho + H/c_{\rm hyp})$ undershoots the true one by
exactly that ratio. This is a spurious sink, strongest exactly where
$c_{\rm hyp}\kappa\rho$ is smallest relative to $H$, i.e.\ at low density
and low metallicity -- the early ultra-faint-dwarf ISM this module is
built for. Verified by a cosmological run of \code{ISRFCosmology}'s
\code{free\_field} fixture, redshift $9$ to $a=0.125$: the corrected form
passes (mass-weighted $u_{\rm FUV}/u_0-1$ and $u_{\rm LW}/u_0-1$ both
$2.2\times10^{-8}$, $900\times$ below their bar), while the bare-$H$
form's own prediction for that run, a decay of order $20\%$, would have
failed catastrophically against it
(\code{.claude/dev/logs/2026-09-18\_1349\_cosmo-hubble-rescale.md}).

The relaxation time is $\tau = \lambda/c_{\rm hyp}$, and the flux equation
is a Cattaneo-type relaxation equation \citep{Cattaneo1958}: in the
isotropic limit its steady state is Fick's law $\mathbf F = -(c_{\rm
hyp}\lambda/3)\,\nabla u$, with an effective diffusion coefficient $D =
c_{\rm hyp}\lambda/3$ that follows from $\lambda$ and $c_{\rm hyp}$
rather than being a free parameter. The reduced flux $f$ and the
constraint $|\mathbf F|\le c_{\rm hyp}u$ are built with $c_{\rm hyp}$, not
$c$: built with $c$, $f$ would be tiny everywhere and the closure would
silently collapse to the isotropic one.

The gas frame is also the only frame in which the scheme is Galilean
invariant. With $c_{\rm hyp}$ no larger than a few times the sound speed,
a box-frame advection term would not be small against the reduced signal
speed, whereas in the gas frame the velocity that remains is the
\emph{relative} one between a source and the gas that absorbs its light,
small for a star moving with its natal cloud. The price is a source moving
at $v_{\rm rel}$ through the gas: it trails a steady field whose centroid
lags it by $v_{\rm rel}\tau$.

\subsubsection{Steady states}
\label{sec:fuv-steady}

Two exact steady states of Eqs.~\ref{eq:fuv-u-eq}--\ref{eq:fuv-F-eq}
for a point source of band luminosity $L$ in a static, uniform medium
of absorption length $\lambda$ ($H=0$) fix what the scheme must
reproduce. In both, $c_{\rm hyp}$ cancels: the reduced speed changes how
long the field takes to settle, not where it settles. In the
free-streaming branch ($f=1$, $\mathbf F_V = c_{\rm hyp} U\,\hat{\mathbf
r}$),
\begin{equation}
  U(r) = \frac{L}{4\pi c\,r^2}\,e^{-r/\lambda} ,
  \label{eq:fuv-steady-freestream}
\end{equation}
the inverse-square law times dust extinction over the physical
absorption length, the exact answer for a point source in a purely
absorbing medium. In the isotropic branch ($f\to0$, $\mathbf D = \mathbf
I/3$), Fick's law turns Eq.~\ref{eq:fuv-u-eq} into a screened Poisson
equation, $\nabla^2 U - (3/\lambda^2)U = -(3L/\lambda c)\,\delta^3(\mathbf
r)$, whose solution is
\begin{equation}
  U(r) = \frac{3L}{4\pi c\,\lambda\,r}\,e^{-\sqrt3\,r/\lambda} ,
  \label{eq:fuv-steady-diffusion}
\end{equation}
with the textbook diffusion screening length $\lambda/\sqrt3$. A
single point source in an absorbing medium never isotropizes (there is
no scattering): the M1 system is solved exactly by the free-streaming
solution, $\mathbf F_V = c_{\rm hyp}\,U\,\hat{\mathbf r}$ ($f=1$), so the closure
selects Eq.~\ref{eq:fuv-steady-freestream} around one star. The
isotropic branch is not a self-consistent steady state that close to a
source: its own reduced flux, obtained from Eq.~\ref{eq:fuv-steady-diffusion}
through Fick's law, exceeds $1$ for $r\lesssim0.79\,\lambda$, so the
admissibility constraint of Eq.~\ref{eq:fuv-admissible} and the flux
limiter of Section~\ref{sec:fuv-limiter} clip it there. It becomes an
available closure only in the outer envelope, and it is the relevant
limit for a diffuse background built from many sources in many
directions, the production use of this module. Consistent with this, the
discrete scheme's own single-source steady state sits on the flux
limiter's boundary, $|\mathbf F_i^\ast| = c_{{\rm hyp},i}u_i$, rather than
in the interior of the admissible set. Integrating Eq.~\ref{eq:fuv-u-eq}
over a periodic domain removes the transport term for any closure and
gives a closure-independent amplitude identity for a steady source,
\begin{equation}
  \sum_i \frac{m_i u_i}{\lambda_i} = \frac{1}{c}\sum_i m_i \dot s_i ,
  \label{eq:fuv-amplitude-identity}
\end{equation}
which the discrete scheme satisfies exactly when every particle shares
the same $h_i$ and $\Delta t_i$, and hence the same $c_{{\rm hyp},i}$ and
relaxation depth $a_i$: only then does the transport operator conserve
$\sum_i m_i u_i$ exactly across the applied update
(Section~\ref{sec:fuv-operators}); under mixed time bins or a
non-uniform $c_{\rm hyp}\kappa\rho$ the identity holds up to the same
residual described there.

\subsubsection{Coupling to the chemistry and cooling solver}
\label{sec:fuv-coupling}

At every cooling call, the two fields are converted to the quantities
Grackle expects. The photoelectric heating strength is the band-summed
energy flux of a quasi-isotropic field, $c\,U$, in Habing units
\citep{Habing1968}, with each band clamped to non-negative before the
sum, not the sum clamped afterwards,
\begin{equation}
  G_{0,i} = \frac{c\,\rho_i\,\big(\max(u_{{\rm FUV},i},0) +
                 \max(u_{{\rm LW},i},0)\big)}
                 {1.6\times10^{-3}\,{\rm erg\,s^{-1}\,cm^{-2}}} .
  \label{eq:fuv-g0}
\end{equation}
Clamping the sum instead would let a negative band, a propagation
undershoot (Section~\ref{sec:fuv-dissipation}), subtract from the other
band rather than contribute nothing: on a fixture with one band negative
and the other positive this delivered $1.50432$ against the unbiased
$2.00576$, both in Habing units, a $25\%$ suppression that raised no
warning because the sum itself stayed positive throughout. Mixed sign
between the two bands is ordinary at source turn-on, not a rare corner
case, so the two clamps cannot be merged into one on the sum.
Grackle \citep{Smith2017} applies it as the heating rate of
\citet[their Eq.~1]{Wolfire1995}, built on the grain photoelectric
mechanism of \citet{BakesTielens1994}, scaled by the
same gas metallicity relative to solar as the dust-to-gas ratio of
Section~\ref{sec:fuv-opacity},
\begin{equation}
  \Gamma_{\rm PE} = A\,\epsilon_{\rm PE}\,G_0\,n_H\,
    \frac{Z}{Z_\odot}
  \qquad (T < 2\times10^4\,{\rm K}) ,
  \label{eq:fuv-pe-heating}
\end{equation}
per unit volume, with $n_H$ the hydrogen nuclei number density
(including H$_2$ nuclei), and scales its dust recombination cooling by
the same $G_0$; H$_2$ formation on dust also responds to $G_0$, but only
indirectly, through the dust temperature $G_0$ helps set, not a direct
scaling like the two terms above.
\code{GrackleCooling:photoelectric\_heating\_efficiency}
selects $A$ and $\epsilon_{\rm PE}$. With \code{constant} (the default)
the efficiency is constant, $A = 10^{-24}\,{\rm erg\,s^{-1}}$ and
$\epsilon_{\rm PE} = 0.05$. With \code{wolfire1995}, $A$ is the same and
$\epsilon_{\rm PE}$ follows \citet[their Eq.~2]{Wolfire1995}, a function
of the grain charging parameter $G_0 T^{1/2}/n_e$, with Grackle's own
electron density. With \code{density\_dependent} (a Grackle build that
provides it) the efficiency is the density-only fit of
\citet[their Eqs.~A1--A2]{Smith2026} to the solar-neighbourhood
predictions of \citet{Wolfire2003},
\begin{equation}
  A = 1.3\times10^{-24}\,{\rm erg\,s^{-1}}, \qquad
  \epsilon_{\rm PE} = \min\!\left[0.07,\;
    0.0149\,\left(\frac{n_H}{{\rm cm^{-3}}}\right)^{0.235}\right] .
  \label{eq:fuv-pe-efficiency-density}
\end{equation}
\citet{Smith2026} write $n$ for the gas number density and normalise the
dust-to-gas ratio to $0.0081$; Grackle uses $n_H$ in both places and
$Z/Z_\odot$, i.e. a dust-to-gas ratio normalised to its own $0.009387$.
Relative to \code{constant}, \code{density\_dependent} heats by a factor
$0.23$, $0.39$, $0.67$, $1.14$ and $1.82$ at $n_H = 0.1$, $1$, $10$, $100$
and $\ge 725\,{\rm cm^{-3}}$. Grackle's documentation flags its electron
density as unreliable in dense, cold gas, which affects
\code{wolfire1995} only.
The H$_2$ photodissociation
rate is a cross-section times the LW photon flux,
\begin{equation}
  k_{{\rm diss},i} = \sigma_{\rm H_2}\,
    \frac{c\,\rho_i\,u_{{\rm LW},i}}{E_{\rm LW}} , \qquad
  \sigma_{\rm H_2} = 2.47\times10^{-18}\,{\rm cm^2}, \qquad
  E_{\rm LW} = 12\,{\rm eV} ,
  \label{eq:fuv-kdiss}
\end{equation}
with $\sigma_{\rm H_2}$ an effective, spectrum-averaged band cross-section,
not an atomic-physics constant. H$_2$ absorbs in the Lyman and Werner
line systems, so the true dissociation rate is a sum over discrete lines
weighted by the radiation field's own shape across $11.2$--$13.6$\,eV;
$\sigma_{\rm H_2}$ is the single number that reproduces that sum for one
assumed shape. \code{RADIATION\_SIGMA\_H2\_LW\_CGS} in
\code{src/feedback/GEAR/radiation.h} holds it as a compile-time constant,
not a run-time parameter, since it is a property of the assumed spectrum
rather than a knob to tune per run.

The assumed shape is the \citet{Draine1978} interstellar field. Both
$\sigma_{\rm H_2}$ and $E_{\rm LW}$ are fixed under that assumption, and
since $k_{\rm diss}\propto\sigma_{\rm H_2}/E_{\rm LW}$ only their ratio is
constrained: the two are calibrated jointly, not independently.
$E_{\rm LW} = 12$\,eV is itself a band-midpoint stand-in, used because a
gas particle carries the LW band's energy flux $u_{{\rm LW},i}$ while
\eqref{eq:fuv-kdiss} needs a photon flux, and the table carries no per-mass
mean LW photon energy to convert with, unlike the ionizing band, whose
mean excess photon energy above $13.6$\,eV is tabulated. A stellar population whose LW
spectrum is much harder or softer than Draine's is therefore represented
only approximately.

The ratio is checked against the free-space photodissociation rate of
\citet[their Eq.~5]{Sternberg2014},
$D_0 = 5.8\times10^{-11}\,I_{\rm UV}\,{\rm s^{-1}}$, where $I_{\rm UV}$
normalises the Draine field integrated over $912$--$1108$\,\AA\ alone.
$D_0$ is the full-$4\pi$ free-space rate, not their slab-surface
$D(0) = D_0/2$, whose geometric factor does not apply here. Converting
conventions with $I_{\rm UV} = G_0/1.7$ on that band, and taking the LW
band to carry a fraction $0.149$ of the combined $6$--$13.6$\,eV Habing
flux that \code{RADIATION\_HABING\_FLUX\_CGS} normalises,
\eqref{eq:fuv-kdiss} gives $k_{\rm diss}/D_0 = 0.898$, within the $\sim
10\%$ uncertainty of the band integrations entering that comparison. Both
expressions are linear in the field strength, so agreement is a fixed
multiplicative offset: the comparison bounds the normalisation and says
nothing about how either behaves under a different spectral shape. The
derivation is reproduced by
\code{theory/GEAR/Radiation/verify\_sigma\_h2\_lw\_sternberg2014.py}, which
reads the constants out of \code{radiation.h} directly so it cannot drift
from the code it checks. Grackle adds
$k_{\rm diss}$ to its background dissociation rate and applies its own
self-shielding factor \citep{WolcottGreenHaiman2019} to the sum. Grackle
adds no heating for these dissociations by default. The optional switch
\code{GrackleCooling:H2\_photodissociation\_heating} (off by default,
a Grackle build that provides it) adds, following
\citet[their Eqs.~37--38]{Kim2023},
\begin{equation}
  \Gamma_{\rm H_2,diss} = \left[0.4 + \frac{2 f_{\rm pump}}
    {1 + n_{\rm crit}/n}\right]{\rm eV}\;k_{\rm diss,shield}\,n_{\rm H_2},
  \qquad f_{\rm pump} = 8 ,
  \label{eq:fuv-h2-diss-heating}
\end{equation}
per unit volume, the sum of their Eq.~37 (UV pumping) and Eq.~38
(dissociation): $0.4$\,eV from the dissociation itself and the
collisionally thermalised part of the UV-pumping energy.
$k_{\rm diss,shield}$ is the shielded total rate (UV background plus
$k_{\rm diss}$). The thermalised fraction $1/(1 + n_{\rm crit}/n)$ uses
the critical density of \citet[their Eq.~23]{Omukai2000} that Grackle
already applies to H$_2$ formation heating,
$n_{\rm crit}/n = n_{\rm cr,n}/(n_{\rm HI}\,n_{\rm cr,d1} + n_{\rm H_2}\,n_{\rm cr,d2})$.
It differs from the critical density of \citet[their Eq.~39]{Kim2023},
$(A_{\rm eff} + k_{\rm diss,shield})/(x_{\rm H}k_{\rm H_2\mbox{-}H} + x_{\rm H_2}k_{\rm H_2\mbox{-}H_2})$
with $A_{\rm eff} = 2\times10^{-7}\,{\rm s^{-1}}$ \citep{Burton1990} and
the de-excitation rates of \citet[their Eq.~A.1]{Visser2018}, by factors
of $0.2$ to $10^5$ over $10$--$10^4$\,K and $0 < x_{\rm H_2} \le 0.5$
(Omukai larger below about $200$\,K, smaller above).
\citet[their Appendix~A]{Smith2026}, which adopts this heating, prints
the reciprocal factor $1/(1 + n/n_{\rm crit})$. Both quantities are clamped to be
non-negative before they reach Grackle; a negative value is an
undershoot of the numerical scheme (Section~\ref{sec:fuv-dissipation})
and is reported. Grackle applies $k_{\rm diss}$ only to the H$_2$
reaction network itself: it adds no heating for the photodissociation
event or the UV pumping that precedes it, unlike the photoelectric and
dust-recombination channels above.

Enabling the module (\code{GEARFeedback:with\_photoelectric\_heating})
configures Grackle for it: dust chemistry on, which selects the
constant-efficiency photoelectric heating and dust recombination cooling
together, and forces H$_2$ formation on dust grains
(\code{h2\_on\_dust}) on in the cooling modes that track H$_2$
(modes 2 and 3); the per-particle ISRF field
instead of the scalar one, for every dust channel; and, in those same
H$_2$-tracking modes, the radiative-transfer rate coupling that carries
$k_{\rm diss}$. Photoelectric heating and dust recombination cooling
work in every Grackle mode; H$_2$ formation on dust and H$_2$
dissociation only where H$_2$ is a tracked species.

\subsection{How we solve them numerically}
\label{sec:fuv-numerics}

\subsubsection{One time step, in order}
\label{sec:fuv-sequence}

Every active gas particle $i$ advances both bands on its own
hydrodynamical time step $\Delta t_i$; there is no sub-cycling and no
separate radiation time step. In order:
\begin{enumerate}
  \item Set this step's signal speed $c_{{\rm hyp},i}$ and absorption
    rates $\kappa_{b}\rho_i$; draw this step's source rate $\dot s_i$
    from the particle's dose reservoir (Section~\ref{sec:fuv-injection}).
  \item Density loop: accumulate the neighbour reference field for the
    dissipation trigger.
  \item Gradient loop: build the closure tensor from $(u^n,\mathbf F^n)$
    and accumulate the pressure-tensor divergence of $u^n$
    (Eq.~\ref{eq:fuv-grad-op}).
  \item Update $\mathbf F_i$ by exact relaxation
    (Eq.~\ref{eq:fuv-F-update}), apply the flux limiter against $u_i^n$
    (Eq.~\ref{eq:fuv-flux-limiter}), and update the dissipation
    coefficients from $u^n$ (Section~\ref{sec:fuv-dissipation}).
  \item Force loop: accumulate the divergence of the new flux
    (Eq.~\ref{eq:fuv-div-op}) and the artificial-dissipation exchange of
    $u^n$ (Eq.~\ref{eq:fuv-diss-formula}).
  \item Update $u_i$ by exact relaxation with the dissipation folded in
    (Eq.~\ref{eq:fuv-u-update}). This is the state cooling reads.
  \item Stars active this step deposit their band luminosities into
    the dose reservoirs of their kernel neighbours
    (Eq.~\ref{eq:fuv-inject}).
\end{enumerate}
Steps 2 to 7 run inside the existing SPH density, gradient and force
loops; the scheme adds pairwise operators to those loops and no new
tasks.

\subsubsection{The signal speed and the stability bound}
\label{sec:fuv-chyp}

The hyperbolic Courant condition for
Eqs.~\ref{eq:fuv-u-eq}--\ref{eq:fuv-F-eq} is $\Delta t \le C\,h/c_{\rm
hyp}$, linear in $h$ like the hydrodynamical one. Rather than derive a
time step from a chosen speed, the scheme keeps the particle's own,
already-decided time step and derives the speed from it,
\begin{equation}
  c_{{\rm hyp},i} = \min\!\left(C_{\rm hyp}\,\frac{h_i}{\Delta t_i},\; c\right) ,
  \label{eq:fuv-chyp}
\end{equation}
so that the Courant condition holds as an equality by construction, per
particle and per step, before the clamp at $c$. The clamp can only make
a step more conservative. $C_{\rm hyp}$ (default $0.5$) is the one
dimensionless dial of the propagation. Its consequence is a transient
that is slow in diffuse gas: at $h\sim10$\,pc and $\Delta t\sim10^5$\,yr,
$c_{\rm hyp}\sim 50\,{\rm km\,s^{-1}}$, and with $\lambda$ of tens of
kiloparsecs at $n_H\sim1\,{\rm cm^{-3}}$ and $Z\sim10^{-2}Z_\odot$, the
relaxation time $\tau = \lambda/c_{\rm hyp}$ exceeds an OB star's
lifetime by a large factor, so such a star's field never reaches its
steady state. This is the cost of the reduced speed, not a defect: at
the true $c$ the same $\tau$ is of order $10^5$\,yr.

The explicit, staggered integrator below is stable only if the
per-step transport number and the per-step dissipation number together
satisfy
\begin{equation}
  6.2\,\alpha\,C_{\rm hyp} + 0.70\,C_{\rm hyp}^2 \le 2 , \qquad
  \alpha = \max(\alpha_{\max}, \alpha_{\rm floor}) ,
  \label{eq:fuv-stability-bound}
\end{equation}
where $\alpha_{\max}$ and $\alpha_{\rm floor}$ are the two ceilings of
the artificial-dissipation coefficient
(Section~\ref{sec:fuv-dissipation}). The two constants are lattice sums
of the Wendland-C2 kernel \citep{Wendland1995,DehnenAly2012} at the
project's neighbour number: $6.2 = 2
I_W$ with $I_W = 3.10$ the dimensionless self-damping sum of the
dissipation operator, and $0.70 = \nu_{\max}^2/2$ with $\nu_{\max} =
1.18\,C_{\rm hyp}$ the largest transport eigenvalue per step. The bound
is checked at start-up. With dissipation off it gives $C_{\rm hyp}\le
1.69$; at the shipped $\alpha = 0.5$ it gives $C_{\rm hyp}\le 0.571$.

\subsubsection{Speed heterogeneity: two failure pathways, an open
residual, and what is settled about each}
\label{sec:fuv-chyp-heterogeneity}

Eq.~\ref{eq:fuv-chyp} is evaluated independently for every particle, so
$c_{{\rm hyp},i}$ generally differs from $c_{{\rm hyp},j}$ for any pair
$(i,j)$ that sits on different time bins or has different $h$ -- a
``seam'' the pairwise transport and injection operators cross on every
step. An investigation opened 2026-09-18 (operator-directed, following
up the T6 validation above) separated this heterogeneity's consequences
into two distinct failure pathways and quantified each against the
shipped scheme's own example fixtures
(\code{examples/SubgridTests/StellarFeedback/ISRF/}). Both are
established; a residual mechanism beyond either is not, and no fix is
adopted here -- \textbf{this section documents an open investigation, not
a resolved design.}

\paragraph{Pathway 1: negative energy at a speed seam.} A receiver
particle's own Courant number is set by its neighbours' $c_{\rm hyp}$,
not only its own; a fast-stepping neighbour with a large $c_{\rm hyp}$
can drive the shared pairwise flux term past what the slower receiver's
own step can absorb, undershooting $u$ past zero. This is now established
as a real cause, three independent ways, on
\code{ISRFDissipationExtreme} and \code{ISRFCausalReach}: pinning every
particle to one uniform $c_{\rm hyp}$ removes the sign violation
entirely (a fixture-specific sign-closure gate goes from $0.57\%$/$0.85\%$
FUV/LW to exactly $0.0000\%$); lowering $C_{\rm hyp}$ alone, with no code
change, removes every negative on a fixture where the artificial
dissipation cannot repair one; and two independently implemented
neighbour-dependent pair weights, predicted in advance to trade this
pathway for a different SPH-consistency violation, both do exactly that,
regressing the same fixture's sign closure by a factor of $10$--$13$. See
the \code{swift-gear} logbook, \code{theory/radiation/}, for the full
numbers and job references; this document states the mechanism, not the
validation record.

\paragraph{Pathway 2: an amplitude error from coupling the source to
$c_{\rm hyp}$.} The source term in Eq.~\ref{eq:fuv-u-eq} is rescaled by
$c_{\rm hyp}/c$ at the injecting particle's own value, so an identical
star deposits a physically different amount of energy depending on the
local gas's own time-stepping: a particle pinned to a short step (large
$c_{\rm hyp}$) receives more of the true source than one on a long step,
independent of the pathway-1 sign question above. At a uniform
$c_{\rm hyp}$ this construction is exactly self-consistent: the source
rescale and the rate getters in \code{radiation\_gas.c} (which read the
true $c$, never $c_{\rm hyp}$) form a matched pair whose propagation
fixed point cancels the speed identically, with or without transport
(same-day adversarial check, \code{.claude/dev/ISRF\_HISTORY.md},
2026-09-18 $\sim$11:55). The defect is specific to $c_{\rm hyp}$
\emph{varying between neighbours}: a particle's divergence sums fluxes
carrying each sender's own $c_{{\rm hyp},j}$ while its own terms carry
$c_{{\rm hyp},i}$, an interface (neighbour-ratio) effect in the sense of
\citet{Katz2017}'s problem class (Section~\ref{sec:fuv-chyp-literature}),
not a general source-normalisation one. At a \emph{uniform} $c_{\rm hyp}$
the cancellation now holds under cosmology too, since $H$ is dilated by
the same $c_{\rm hyp}/c$ factor as every other rate
(Section~\ref{sec:fuv-lagrangian}, \code{d47419d04}); an earlier version
of this paragraph called this unmeasured because that dilation was
missing at the time. Whether the cancellation still holds when
$c_{{\rm hyp},i}$ varies \emph{between neighbours}, the actual subject of
this pathway, is unmeasured under cosmology and needs an operator ruling
before a cosmological run with heterogeneous speeds is trusted on this
point.

What is measured: forcing a uniform speed (the pin, and independently a
lowered $C_{\rm hyp}$) changes the final FUV band energy on
\code{ISRFCausalReach} relative to the shipped, heterogeneous-speed run;
the fixture is still rising at $t_{\rm end}$ under every speed tested, and
the same-day check above found this fixture's own cross-arm ratio has not
settled by that point (still declining across snapshots), so no number
from it is quoted here. Part of any gap measured this way is a slower
relaxation clock (Section~\ref{sec:fuv-chyp}'s own discussion of the
transient), not a pure injection bias, and the two effects are not
separated on the fixtures tested so far.

\paragraph{Speed definitions tested, and the default chosen among them.}
Four alternative definitions of $c_{{\rm hyp},i}$ and/or the pairwise
operators were implemented as selectable runtime schemes
(\texttt{GEARFeedback:ISRF\_c\_hyp\_scheme}), against scheme $0$ (the
formula above, unmodified operators) as baseline: a kernel-local speed
rule that takes the slowest step among a particle and every neighbour in
its kernel (scheme $1$); a fixed global fraction of $c$, independent of
$h$ or the time step (scheme $2$); a ``consistent variable-$c$''
reformulation that changes which variable the pairwise operators carry so
that every operator becomes the per-particle-speed generalisation of the
reduced-speed-of-light method, reducing bit-for-bit to the shipped scheme
under a uniform speed (scheme $3$,
Section~\ref{sec:fuv-variable-c-operators}); and the combination of the
kernel-local speed with that operator rewrite (scheme $4$). All four
measurably reduce pathway 1 relative to shipped on every fixture tested;
scheme $4$ additionally clears most of Pathway 2's amplitude defect,
which schemes $1$ and $2$ alone do not touch. None of $1$, $2$, $3$ alone
clears its own pre-registered acceptance bar: each fails at least one
fixture's exact sign-closure or accuracy criterion; the fixed-fraction
scheme costs nothing extra at a fraction low enough to be free, but only
because that fraction is below the bulk sound speed, making the
radiation non-causal, while at the only fraction so far judged physically
defensible it carries an initially undercounted particle-update cost of
$+21$ to $+29\%$ once measured correctly (a low fraction can never bind
the new timestep term, so it was never informative about cost); and the
kernel-local scheme's own neighbour-maximum accumulates only over
neighbours within $i$'s own kernel, so a larger-$h$ receiver that is
inside a sender's kernel but has the sender outside its own reads that
sender's flux without ever including it in its speed maximum -- a
dense-short-step sender beside a diffuse-long-step receiver, which a
uniform-resolution test fixture cannot exercise but a real galaxy's
density contrast can, so this rule reduces pathway 1 without removing it
``by construction''.

\textbf{Scheme $4$ is the shipped default}, as of \code{a53ab8311}: of
the five, it has the lowest negativity in $7$ of $8$ fixture-band pairs
measured (two of them exact zeros) and the best amplitude,
$1.088\times$ the exact fixed point against the shipped scheme's
$3.06\times$, at bit-identical particle-update cost (neither the
kernel-local speed formula nor the operator rewrite alone changes what
work is done per particle). Scheme $0$ remains reachable by setting
\code{ISRF\_c\_hyp\_scheme} explicitly. \textbf{This ranking rests on
three fixtures only}, which the operator's own acceptance criterion for
a default does not by itself consider sufficient
(\code{.claude/dev/ISRF\_MASTER.md}); a validation campaign across every
ISRF example, at several resolutions and including a cosmological run,
is prepared and partly submitted as of this writing but not complete.
\textbf{The default may yet be reverted to scheme $0$} if that campaign
fails; this section reports the current shipped choice and its
measured justification, not a closed validation.

\paragraph{Open question: is pathway 1 the whole story?} A same-day
finding (2026-09-18, not yet independently re-tested) proposes that the
scheme's central-SPH-divergence-plus-pairwise-dissipation structure is
positive only in the diffusive regime and is \emph{not} positive for a
free-streaming neighbour (M1 flux ratio $f\to1$) at any choice of
$c_{\rm hyp}$, uniform or not, once the required dissipation coefficient
exceeds what the shipped stability bound (Eq.~\ref{eq:fuv-stability-bound})
allows -- which would mean the speed-seam heterogeneity mainly amplifies
a pre-existing structural gap rather than being its root cause. This is
a working hypothesis under active test as of this writing, not an
established result: it is included here so the open question is stated
precisely, not so a reader treats it as settled. It must not be read as
a second candidate fix alongside the three above.

\subsubsection{Literature position on a variable propagation speed}
\label{sec:fuv-chyp-literature}

A literature review conducted 2026-09-18
(\code{.claude/dev/design-isrf-rsla-literature.md} in this checkout)
answered whether any published radiation-transport scheme varies its
reduced speed of light the way Eq.~\ref{eq:fuv-chyp} does. Citations
below are separated as the review itself separates them: \textbf{verified}
means the source's full text was read directly (or, for the two SWIFT
modules, the source code); \textbf{secondary} means only a search
summary or a citing paper's own reference list was checked, and no
number from such a source should be trusted without re-reading it.

\textbf{No published scheme ties the reduced speed of light to a
per-element time step (verified).} Where the literature does vary the
speed spatially, it is tied to a static hierarchy -- the AMR refinement
level -- with an exact, fixed ratio (a factor $2$) between any two
adjacent levels, not to an independently adaptive per-element quantity
\citep{Katz2017,Rosdahl2018}. \citet{Katz2017} give an exact interface
treatment for this case (their Appendix A): the neighbour's photon
density is rescaled by the ratio of the two levels' speeds before
entering a shared Lax-Friedrichs flux, which both sides of the interface
compute identically, giving flux-exact conservation for a single common
time step across all levels. This construction does not port to GEAR's
SPH scheme: it relies on exactly two well-defined neighbour states per
face and one shared flux value both sides agree on, whereas SPH's
pairwise transport has each particle accumulate its own side of an
interaction independently, with no shared face. Even in its native
finite-volume setting, the exact proof degrades to an admitted
few-percent, not sign-violating, leak once per-level time steps are
allowed to desynchronise \citep{Katz2017,Rosdahl2018}. A citation inside
\citet{Katz2017} (Commer\c{c}on, Debout \& Teyssier 2014, A\&A 563, A11)
is flagged by the review as the one lead that addresses this
asynchronous-time-step case directly; it was not independently read and
its content is not asserted here.

\textbf{SWIFT's own two RT modules both use a single, global, run-constant
value (verified, read directly from \code{src/rt/}).} GEAR-RT
(\code{src/rt/GEAR/rt\_properties.h}) sets \code{reduced\_speed\_of\_light}
once at start-up from the parameter \code{GEARRT:f\_reduce\_c} and every
particle reads the same global struct field. SPHM1RT
(\code{src/rt/SPHM1RT/rt\_properties.h}) does the same with
\code{SPHM1RT:cred}, copied unchanged into every particle at start-up. A
single uniform speed is not a deviation from how this feature is
normally used elsewhere in this codebase; every other user of a reduced
speed of light in \swift{} uses exactly one number for the whole run.

\textbf{The source/injection amplitude must not be coupled to the
reduced speed at all (verified).} \citet{Gnedin2026} states this
directly: ``this is the \emph{only} place where the speed of light
should be reduced[;] \dots the speed of light should \emph{not} be
reduced in the equation for the cosmic background\dots [n]or should $c$
that enters the definition of $\psi_\nu$ \dots be reduced, as it would
result in the incorrect photon production rate.'' \citet{Ocvirk2019}
confirm this numerically in an independent code (RAMSES-CUDATON):
reducing $c$ by a factor $f$ inflates the post-reionisation neutral
fraction by the same factor $f$ when the reduced speed leaks into the
chemistry coupling, and the error disappears when the true $c$ is used
there instead. Both papers warn against exactly this operation when the
reduction is \emph{unmatched}: a source rescaled by $c_{\rm hyp}/c$ whose
rate is then read back at the reduced speed. GEAR's convention pairs the
source rescale with a rate getter that reads the true $c$
(Section~\ref{sec:fuv-chyp}, Pathway 2), which a same-day check found to
cancel exactly at a uniform speed; the residual defect these citations
bound is therefore the per-neighbour interface case (heterogeneous
$c_{\rm hyp}$, Pathway 2 above), not a general defect at any single
global speed.

\textbf{The RSLA validity criterion in the literature is derived for an
ionizing, recombination-bounded front and does not transfer directly to
GEAR's non-ionizing bands (verified for the cited numbers, non-transfer
is this review's own reading).} \citet{Rosdahl2013}'s criterion
($f_c\gtrsim10\,t_{\rm cross}/\tau_{\rm sim}$) and the values people
actually run at (STARFORGE, \citealt{Grudic2021}, single global
$\tilde c=5$--$300\,{\rm km\,s^{-1}}$, fiducial $90\,{\rm km\,s^{-1}}$)
are both set from ionizing-front or recombination timescales that have
no analogue for photoelectric heating or H$_2$ photodissociation. No
source found by the review gives a validity criterion or a value
specific to non-ionizing ISM radiation transport; the general form that
does transfer is simply that $c_{\rm hyp}$ must exceed every physical
signal speed the transport is meant to resolve, evaluated against the
shortest timescale of interest.

\textbf{No prior report of negative radiation energy at a variable-speed
interface was found (verified absence, not a positive claim).} The
review checked every RSLA-interface source it read directly, including
the paper most focused on RSLA conservation failures generally
\citep{Gnedin2026}, for any mention of a sign-violating failure at a
speed boundary; none was found. This may mean the failure is rare in
schemes that never couple a per-element time step to the propagation
speed the way this scheme does, or it may mean it is undocumented; the
review could not distinguish the two, and the unread Commer\c{c}on et
al. (2014) source above is the most promising remaining lead.

\subsubsection{Injection: kernel deposit, extinction and dose reservoir}
\label{sec:fuv-injection}

A star $\star$ active on a step of length $\Delta t_\star$ deposits, on
each gas particle $j$ inside its kernel, the band energy
\begin{equation}
  \Delta E_{b,j} = \Delta t_\star\, w_{j\star}\, L_{b,\star}\,
    \zeta_{b,j} , \qquad
  w_{j\star} = \frac{m_j\,W(r_{j\star}, h_\star)}
                    {\sum_k m_k\,W(r_{k\star}, h_\star)} ,
  \label{eq:fuv-inject}
\end{equation}
with weights that sum to one over the kernel. $\zeta_{b,j}\in(0,1]$ is
a receiver-side dust extinction,
\begin{equation}
  \zeta_{b,j} = \exp\!\big(-\kappa_b(Z_j)\,\Sigma_j\big) , \qquad
  \Sigma_j = 2\,H_j\,\rho_j ,
  \label{eq:fuv-extinction}
\end{equation}
where $H_j$ is the kernel support radius of particle $j$: the column
through the receiving particle's own resolution element. It represents
dust immediately local to the receiver (a particle inside its own dense
clump) and is distinct from the absorption along the path, which the
propagation handles. Extinction applies in both run modes.
$\Sigma_j = 2\,H_j\,\rho_j$ is the default (\texttt{GEARFeedback:%
ISRF\_extinction\_path} = \texttt{kernel\_diameter}); setting it to
\texttt{kernel\_radius} instead uses $H_j\,\rho_j$.

At \texttt{COOLING\_GRACKLE\_MODE} $>1$, Grackle's own H2
self-shielding needs a separate column length, set independently via
\texttt{GrackleCooling:H2\_self\_shielding\_path}. At mode 2
(\texttt{H2\_self\_shielding} = 2, kernel support radius), Grackle forms
$N_{\rm H2} = 2 n_{\rm H2} l$, so the length passed is
$l = {\rm path}/2$ with ${\rm path} = {\rm path\_in\_kernel\_radii}\,
\gamma_K h$ (physical); at mode 3 the length is Grackle's own local
Jeans length instead. Dust extinction and H2 self-shielding use
different lengths on purpose: the dust column is local to the
receiver's own resolution element, while H2 self-shielding needs the
scale over which the H2 column actually forms (the Jeans length at mode
3); collapsing both onto one common length would cost accuracy in
whichever regime it is not tuned for.

\textbf{Dose reservoir.} A gas particle usually sits on a finer time
bin than its star. Depositing the whole star-step energy in one lump
would produce a sawtooth in $u_j$ between star visits, with cooling
reading the trough. Instead the deposit goes into a per-particle,
per-band reservoir $R_{b,j}$ (an energy per unit gas mass), and the star's visit
extends the reservoir's drain horizon $t_{{\rm end},j}$ to at least one
star step past the visit:
\begin{equation}
  R_{b,j} \mathrel{+}= \frac{\Delta E_{b,j}}{m_j} , \qquad
  t_{{\rm end},j} \leftarrow \max\!\big(t_{{\rm end},j},\; t + \Delta t_\star\big) .
  \label{eq:fuv-reservoir}
\end{equation}
Any number of stars on any time bins add into the same reservoir. At
the start of each of its own steps the particle drains a fraction of
the reservoir into a source rate that is constant over the step,
\begin{equation}
  \phi_j = \min\!\left(1,\;\frac{\Delta t_j}{t_{{\rm end},j} - t}\right) ,
  \qquad
  \dot s_{b,j} = \frac{\phi_j\,R_{b,j}}{\Delta t_j} , \qquad
  R_{b,j} \leftarrow (1-\phi_j)\,R_{b,j} ,
  \label{eq:fuv-drain}
\end{equation}
which empties the reservoir exactly at the horizon and keeps the rate
constant across the particle's sub-steps between visits. The rate
enters the relaxation update of $u$ as a frozen source, rescaled by
$c_{\rm hyp}/c$ there (Eq.~\ref{eq:fuv-u-update}); the reservoir itself
holds the unrescaled physical dose. If a particle's neighbour search
fails on a step, the drawn amount is returned to the reservoir. The
reservoir is specific to the particle mass at deposit time and is not
rescaled if that mass later changes.

\textbf{Illumination window.} A visit also marks the particle as
illuminated for two star steps and, on the first visit only,
synchronises the particle's time step so that it integrates the deposit
promptly. The mark expires once no star has renewed it for two star
steps; with propagation on the expiry has no effect on the field, which
decays through its own equations.

\subsubsection{Spatial operators}
\label{sec:fuv-operators}

Both transport operators are plain pairwise SPH kernel-gradient
estimators on the particle positions, with no matrix correction, no
face reconstruction and no Riemann solver. Write $\mathbf x_{ij} =
\mathbf x_i - \mathbf x_j$, $r_{ij} = |\mathbf x_{ij}|$, and
$\nabla_i W_{ij}^{(i)} = W'(r_{ij}/h_i)\,h_i^{-4}\,\mathbf x_{ij}/r_{ij}$
for particle $i$'s own kernel gradient at its own smoothing length
(likewise $\nabla_i W_{ij}^{(j)}$ at $h_j$). All densities below are the
step's frozen values.

\textbf{Flux divergence.} A single shared pair coefficient, built from
both particles' own terms, credited to one side and debited from the
other:
\begin{equation}
  \Phi_{ij} = \frac{\mathbf F_i\cdot\nabla_i W_{ij}^{(i)}}{\rho_i}
            + \frac{\mathbf F_j\cdot\nabla_i W_{ij}^{(j)}}{\rho_j} ,
  \qquad
  \Big(\tfrac{1}{\rho}\nabla\cdot(\rho\mathbf F)\Big)_i
    = \sum_j m_j\,\Phi_{ij} ,
  \label{eq:fuv-div-op}
\end{equation}
with the mirrored contribution $-m_i\,\Phi_{ij}$ added to particle
$j$'s own sum. The mirrored credit and debit make $\sum_i m_i\,(\nabla\cdot)_i = 0$
identically for any $h_i\ne h_j$, $\rho_i\ne\rho_j$, provided both sides
of every pair are visited. The density and gradient loops reach particle
$i$ only for $r_{ij} < H_i$, so a pair with $H_i \le r_{ij} < H_j$ would
credit $j$ and never debit $i$; the divergence is therefore accumulated
in the force loop, which visits a pair whenever either kernel reaches.
The gradient operator below needs no such care: $i$'s own term carries
$\nabla_i W_{ij}^{(i)}$, which vanishes for $r_{ij}\ge H_i$. What
each particle applies, however, is that divergence scaled by its own
$\Delta t_i\,\varphi(a_i)$ in Eq.~\ref{eq:fuv-u-update}, and $a_i$
differs between particles whose time bin or $c_{\rm hyp}\kappa\rho$
differ. $\sum_i m_i u_i$ is therefore conserved exactly only when both
particles of a pair share the time bin and the same $c_{\rm
hyp}\kappa\rho$; otherwise a residual of order the difference between
the two particles' $\Delta t\,\varphi(a)$ remains, the price of applying
an exact exponential integrator per particle rather than one global
step. An active/inactive pair, which leaves one side undebited until
that particle's own next step, is a special case of the same effect.

\textbf{Pressure-tensor divergence.} A difference estimator on the
volumetric field $U = \rho u$ contracted with each particle's own
closure tensor, with each particle's own kernel derivative and a
$1/\rho^2$ normalisation:
\begin{equation}
  \mathbf T_{ij} = U_i\,\mathbf D_i\cdot\hat{\mathbf x}_{ij}
                 - U_j\,\mathbf D_j\cdot\hat{\mathbf x}_{ij} , \qquad
  \mathbf G_i \equiv
  \Big(\tfrac{1}{\rho}\nabla\cdot(\mathbf D\,\rho u)\Big)_i
    = \sum_j \frac{m_j}{\rho_i^2}\,
      \mathbf T_{ij}\,\big|\nabla_i W_{ij}^{(i)}\big| ,
  \label{eq:fuv-grad-op}
\end{equation}
with $\hat{\mathbf x}_{ij} = \mathbf x_{ij}/r_{ij}$ and the mirrored
term for $j$ using $\rho_j$ and $\nabla_i W_{ij}^{(j)}$. $\mathbf D_i$ is
recomputed from the particle's own $(u_i^n, \mathbf F_i^n, c_{{\rm
hyp},i})$ through Eq.~\ref{eq:fuv-m1-closure} in every pair interaction
that needs it, rather than cached once per particle; it is never stored,
and every recomputation over one step draws on the same frozen inputs.
The reduced flux is clamped at $1$ because $\mathbf F_i$
was limited against the previous step's $u_i$ (Section~\ref{sec:fuv-limiter}),
and $f = 0$, $\hat{\mathbf n} = \mathbf 0$ are used where $u_i\le0$ or
$\mathbf F_i = \mathbf 0$. At $\mathbf D = \mathbf I/3$ the operator
reduces to $\tfrac13$ of the ordinary SPH estimate of
$\rho^{-1}\nabla(\rho u)$.

The two operators are chosen as a pair. With the shared-coefficient
divergence and the own-derivative, $1/\rho^2$ gradient, the discrete
transport is exactly skew-adjoint in the inner product $\langle a, b
\rangle = \sum_i m_i\rho_i\,a_i b_i$ weighted by $\mathbf D^{-1}$
whenever neighbouring tensors agree ($\mathbf D_i = \mathbf D_j$, which
includes the $\mathbf F = \mathbf 0$ far field). This is the discrete
image of the quadratic invariant $\int (U^2 + \mathbf F_V\cdot\mathbf
D^{-1}\mathbf F_V/c_{\rm hyp}^2)$ that the continuum M1 transport
conserves for frozen $\mathbf D$; where $\mathbf D$ varies between
neighbours an $O(h|\nabla\mathbf D|)$ consistency error remains, and in
the free-streaming limit the metric degenerates ($\mathbf D$'s
transverse eigenvalue vanishes), so stability there rests on the Courant
bound and on the artificial dissipation, not on this property. Whether
the staggered integrator's stability (Section~\ref{sec:fuv-time}) needs
more than this weaker, $\mathbf D_i=\mathbf D_j$-only adjointness is not
resolved here.

\subsubsection{The consistent variable-speed operators (schemes 3 and
4)}
\label{sec:fuv-variable-c-operators}

Eqs.~\ref{eq:fuv-div-op} and \ref{eq:fuv-diss-formula} are exact only for
a single, uniform $c_{\rm hyp}$: the reduced-speed-of-light method is a
rescaling of every rate in
Eqs.~\ref{eq:fuv-u-eq}--\ref{eq:fuv-F-eq} by $c_{\rm hyp}/c$, and that
rescaling commutes with the pairwise sum only when the same $c_{\rm hyp}$
multiplies every term. Its correct generalisation to a per-particle
$c_{{\rm hyp},i}$ (Section~\ref{sec:fuv-chyp-heterogeneity}) is a change
of variable, not a new pair weight, so it leaves the SPH kernel-sum
identity untouched -- unlike the neighbour-dependent pair weights tried
and rejected in that section. Writing the reduced flux
$\tilde{\mathbf F}_i = \mathbf F_i/c_{{\rm hyp},i}$ (so the admissibility
constraint, Eq.~\ref{eq:fuv-admissible}, becomes $|\tilde{\mathbf
F}_i|\le u_i$),
\begin{equation}
  \frac{du_i}{dt} = c_{{\rm hyp},i}\left(-\kappa\rho_i u_i
    - \big(\nabla\cdot\tilde{\mathbf F}\big)_i + \frac{\dot s_i}{c}\right) ,
  \qquad
  \frac{d\tilde{\mathbf F}_i}{dt} = c_{{\rm hyp},i}\Big(-\kappa\rho_i
    \tilde{\mathbf F}_i - \mathbf G_i[u]\Big) .
  \label{eq:fuv-variable-c-eqs}
\end{equation}
Three things change relative to Section~\ref{sec:fuv-operators}.
\textbf{(1)} \code{specific\_flux} stores $\tilde{\mathbf F}$ rather than
$\mathbf F$, so a time-bin change rescales nothing and the
temporal-impedance jump at a speed seam (Pathway 1,
Section~\ref{sec:fuv-chyp-heterogeneity}) disappears. \textbf{(2)} The
divergence's shared coefficient $\Phi_{ij}$ of Eq.~\ref{eq:fuv-div-op},
already owner-normalised once $\tilde{\mathbf F}$ is the stored state, is
multiplied by the \emph{receiver's} own speed rather than the sender's:
particle $i$ receives $c_{{\rm hyp},i}\,m_j\,\Phi_{ij}$, particle $j$
receives $-c_{{\rm hyp},j}\,m_i\,\Phi_{ij}$. \textbf{(3)} The
dissipation's signal speed of Eq.~\ref{eq:fuv-diss-formula} becomes two
receiver-side speeds, $\alpha_{ij}c_{{\rm hyp},i}$ and $\alpha_{ij}c_{{\rm
hyp},j}$, rather than the shared $\alpha_{ij}\min(c_{{\rm hyp},i},c_{{\rm
hyp},j})$.

Both changes make the pair's two sides carry different scalars, so the
plain mass-weighted sum $m_i X_i + m_j X_j$ that
Eqs.~\ref{eq:fuv-div-op} and \ref{eq:fuv-diss-formula} conserve exactly
under a uniform speed is no longer conserved. Dividing each side by its
own $c_{\rm hyp}$ restores it,
\begin{equation}
  \frac{m_i X_i}{c_{{\rm hyp},i}} + \frac{m_j X_j}{c_{{\rm hyp},j}} = 0 ,
  \label{eq:fuv-variable-c-conservation}
\end{equation}
because both changes above reduce to ``multiply by the receiver's own
$c_{\rm hyp}$'' applied to what was otherwise the same shared,
antisymmetric-under-$i\leftrightarrow j$ coefficient; checked directly by
\code{tests/testRadiationISRFForceDispatchConservation.c}. The M1 closure
tensor (Eq.~\ref{eq:fuv-m1-closure}) is built with $c_M=1$ under this
scheme, since $|\tilde{\mathbf F}_i|/u_i = |\mathbf F_i|/(c_{{\rm
hyp},i}u_i)$ identically, the same physical reduced flux $f$ the shipped
scheme computes from $(\mathbf F_i, c_{{\rm hyp},i})$ directly. Every
rewritten operator reduces bit-for-bit to its shipped counterpart
whenever $c_{{\rm hyp},i}$ is spatially uniform, under a build that
disables floating-point reassociation across the shared coefficient;
this bit-identity is verified under \code{clang} only, since \code{GCC}
has no block-scoped equivalent of the pragma that disables it, and does
not carry over to a \code{GCC} (cluster) build.

Scheme $3$ keeps $c_{{\rm hyp},i}$ at Eq.~\ref{eq:fuv-chyp}'s own formula
and applies only this operator rewrite; scheme $4$ additionally uses the
kernel-local speed of Section~\ref{sec:fuv-chyp-heterogeneity}'s scheme
$1$. Scheme $4$ is the shipped default as of \code{a53ab8311}; scheme $0$
(Eqs.~\ref{eq:fuv-div-op}--\ref{eq:fuv-diss-formula} unmodified) remains
reachable by setting \code{GEARFeedback:ISRF\_c\_hyp\_scheme} explicitly
(Section~\ref{sec:fuv-chyp-heterogeneity} gives the full ranking and the
default's caveats).

\subsubsection{Time integration: exact relaxation, staggered}
\label{sec:fuv-time}

Both moment equations are linear decays with the same rate plus a
transport term that is frozen over the step. Each is integrated exactly
for frozen right-hand side. With the relaxation depth and integrating
factor
\begin{equation}
  a_i = c_{{\rm hyp},i}\left(\kappa\rho_i + \frac{H}{c}\right)\Delta t_i , \qquad
  \varphi(a) = \frac{1-e^{-a}}{a}
  \quad (\varphi\to1 \text{ as } a\to0,\ \ \varphi\to1/a \text{ as } a\to\infty),
  \label{eq:fuv-phi}
\end{equation}
$H$ entering divided by the true speed of light $c$, not
$c_{{\rm hyp},i}$ (dilated by the same $c_{\rm hyp}/c$ factor as
absorption, as Section~\ref{sec:fuv-lagrangian} derives), the flux
update, from the pressure-tensor divergence of $u^n$, reads
\begin{equation}
  \mathbf F_i^{\,n+1} = e^{-a_i}\,\mathbf F_i^{\,n}
    - c_{{\rm hyp},i}^2\,\Delta t_i\,\varphi(a_i)\,\mathbf G_i[u^n] ,
  \label{eq:fuv-F-update}
\end{equation}
and the energy update, from the divergence of the \emph{newly} updated
fluxes,
\begin{equation}
  u_i^{\,n+1} = e^{-a_i}\Big[u_i^{\,n}
      + \Delta t_i\,\varphi(a_i)\,\dot u_i\big|_{\rm diss}\Big]
    + \Delta t_i\,\varphi(a_i)\left[\frac{c_{{\rm hyp},i}}{c}\,\dot s_i
      - \Big(\tfrac{1}{\rho}\nabla\cdot(\rho\mathbf F^{\,n+1})\Big)_i\right] ,
  \label{eq:fuv-u-update}
\end{equation}
with the artificial dissipation of Section~\ref{sec:fuv-dissipation}
evaluated on $u^n$. This is a staggered, leapfrog-like scheme: $\mathbf F$
first from the old $u$, then $u$ from the new $\mathbf F$. Its
amplification matrix has the same trace and determinant as the opposite
order ($u$ first), so the stability condition below is unchanged
(\texttt{verify\_isrf\_timestepping\_stability.py}, S2'). The order is
forced by conservation: only the force loop visits both sides of every
pair, and it runs after the flux update. The exponential
integrator is unconditionally stable in the relaxation for any $\Delta
t/\tau$, which matters because $\tau$ is far shorter than the hydro
step in dense gas ($\tau\sim700$\,yr at $c_{\rm hyp}\sim50\,{\rm
km\,s^{-1}}$ in a $10^4\,{\rm cm^{-3}}$ clump at solar metallicity) and
far longer in diffuse gas; a forward-Euler treatment would need
sub-cycling wherever $\Delta t >
\tau$. The
exact solution of the inhomogeneous decay, $u(\Delta t) = u_0 e^{-a} +
\dot s\,\tau\,(1-e^{-a})$, is what Eq.~\ref{eq:fuv-u-update} realises
for the source, which is why injection is folded into the same update
rather than added afterwards.

The flux update has an exact fixed point for a frozen gradient,
\begin{equation}
  \mathbf F_i^\ast = -\,\mathcal C_i\,\mathbf G_i , \qquad
  \mathcal C_i = \frac{c_{{\rm hyp},i}}{\kappa\rho_i + H/c} ,
  \label{eq:fuv-fixed-point}
\end{equation}
which is the discrete form of Fick's law with $D = c_{\rm hyp}\lambda$
in front of the tensor-weighted gradient at $H=0$ (the $1/3$ sits inside
$\mathbf D$); $H/c$ is the same dilated Hubble contribution as in $a_i$
above. The dissipation floor gate of
Section~\ref{sec:fuv-dissipation} measures each particle's distance
from this fixed point.

\subsubsection{The flux limiter}
\label{sec:fuv-limiter}

Once per step, right after the flux update, the flux is rescaled onto
the admissible set against $u_i^n$, the field its closure tensor was
built from:
\begin{equation}
  \mathbf F_i \leftarrow \mathbf F_i\,
    \min\!\left(1,\ \frac{c_{{\rm hyp},i}\,u_i}{|\mathbf F_i|}\right)
    \quad (u_i>0) , \qquad
  \mathbf F_i \leftarrow \mathbf 0 \quad (u_i\le0) .
  \label{eq:fuv-flux-limiter}
\end{equation}
The limiter belongs to the M1 closure, not to the artificial dissipation
of Section~\ref{sec:fuv-dissipation}. The Eddington factor of
Eq.~\ref{eq:fuv-m1-closure} contains $\sqrt{4-3f^2}$ and is derived for,
and physically meaningful on, the admissible range $f\in[0,1]$ alone;
beyond $f = 1$ it describes nothing, and the radical itself turns negative
at $f = 2/\sqrt3$. The flux update of Eq.~\ref{eq:fuv-F-update} is
explicit in the pressure-tensor divergence and can therefore overshoot
$|\mathbf F| = c_{\rm hyp}u$ wherever $u$ is small while $\mathbf F$ is
not, which is exactly the situation at a front. The limiter restores
admissibility once per step, so that the state the next step's closure is
built from satisfies Eq.~\ref{eq:fuv-admissible} by construction; the
$\min(f,1)$ inside the closure tensor (Eq.~\ref{eq:fuv-grad-op}) is a
guard behind it, which with the limiter in place rarely binds. The limiter
touches only $\mathbf F$ and conserves $\sum_i m_i u_i$.

\subsubsection{Artificial dissipation}
\label{sec:fuv-dissipation}

The pairwise kernel-gradient transport has no built-in numerical
dissipation, unlike a Riemann-solver flux, so a steep gradient next to
a strong source or a sharp density contrast can drive $u$ negative.
The remedy is the same as artificial viscosity in SPH hydrodynamics: a
numerical conductivity that acts on the field
\citep{ClearyMonaghan1999,Price2008}, is switched on where it is needed,
and represents no physical process. It must not be confused
with the physical absorption term of Eq.~\ref{eq:fuv-u-eq}.

\textbf{The term.} Per pair and per band, with $U = \rho u^n$ the
volumetric form of the step's starting field and $\bar W'_{ij} =
\tfrac12\big(|\nabla_i W_{ij}^{(i)}| + |\nabla_i W_{ij}^{(j)}|\big)$
the symmetrised radial kernel derivative,
\begin{equation}
\begin{aligned}
  v_{{\rm sig},ij} &= \alpha_{ij}\,\min\!\big(c_{{\rm hyp},i},\,c_{{\rm hyp},j}\big) ,
  \qquad
  \Psi_{ij} = -\,v_{{\rm sig},ij}\,\frac{U_i - U_j}{\rho_i\rho_j}\,\bar W'_{ij} , \\
  \dot u_i\big|_{\rm diss} &= \sum_j m_j\,\Psi_{ij} ,
\end{aligned}
  \label{eq:fuv-diss-formula}
\end{equation}
with the mirrored contribution $-m_i\,\Psi_{ij}$ added to particle
$j$'s own sum. The exchange itself is antisymmetric between the pair,
but it is applied through the same $e^{-a_i}\Delta t_i\,\varphi(a_i)$
factor inside the relaxation update (Eq.~\ref{eq:fuv-u-update}), so $\sum_i m_i u_i$
is conserved exactly over a pair only when both particles share the time
bin and $c_{\rm hyp}\kappa\rho$, for the same reason as the flux
divergence operator above. In the continuum
limit it is a diffusion of $U$ with coefficient of order $\alpha\,v_{\rm
sig}\,h$. The signal speed uses the smaller of the two particles'
$c_{\rm hyp}$: because $c_{{\rm hyp},i}\propto h_i/\Delta t_i$, a
particle pinned to a short time step has a $c_{\rm hyp}$ far above its
neighbours' for numerical rather than physical reasons, and $\min$ is
the combination that keeps the per-step dissipation number bounded
independently of the time-bin ratio across the pair.

\textbf{The pair coefficient} combines two per-particle, per-band
components, a reactive trigger and an anticipatory floor,
\begin{equation}
  \alpha_{ij} = \max\!\big(\alpha_{{\rm trig},i},\ \alpha_{{\rm trig},j},\
                            \alpha_{{\rm floor},i},\ \alpha_{{\rm floor},j}\big) ,
  \label{eq:fuv-diss-alpha-pair}
\end{equation}
so whichever particle needs dissipation gets it at full strength with
every neighbour.

\textbf{The negativity trigger.} The reactive component rises only
when a particle's own field has gone negative relative to its
neighbours' field scale:
\begin{equation}
  \varepsilon_i = \frac{\max(-U_i, 0)}{\max\!\big(\langle|U|\rangle_i,\,-U_i\big)}
  \in[0,1], \qquad
  x_i = \min\!\Big(\frac{\varepsilon_i}{\varepsilon_1}, 1\Big), \qquad
  \alpha_{{\rm aim},i} = \alpha_{\max}\,x_i^2\,(3-2x_i) ,
  \label{eq:fuv-diss-trigger}
\end{equation}
with $\langle|U|\rangle_i = \sum_j (m_j/\rho_j)\,W_{ij}\,|U_j|$ the
kernel mean of the neighbours' previous-step field and $\varepsilon_1$
(default $0.01$) the relative undershoot at which the smooth ramp
saturates at $\alpha_{\max}$ (default $0.5$). The coefficient jumps to
$\alpha_{\rm aim}$ when that is higher than its current value and
otherwise decays towards it, the switch form of \citet{CullenDehnen2010}
rather than the source-and-decay equation of \citet{MorrisMonaghan1997},
\begin{equation}
  \alpha_{{\rm trig},i} \leftarrow \alpha_{{\rm aim},i}
    + \big(\alpha_{{\rm trig},i} - \alpha_{{\rm aim},i}\big)
      \exp\!\left(-\frac{c_{{\rm hyp},i}\,\Delta t_i}{5\,h_i}
                  - c_{{\rm hyp},i}\,\kappa\rho_i\,\Delta t_i\right) ,
  \label{eq:fuv-diss-decay}
\end{equation}
an $e$-folding over five smoothing lengths of signal travel, faster in
optically thick gas. There is no velocity-divergence switch because the
field carries no bulk velocity; the switch is built from the field's own
sign. The two bands carry independent coefficients.

\textbf{The floor.} The trigger is exactly zero on the positive front
of a pulse and therefore cannot damp the dispersive wake that the front
sheds behind it in the optically thin regime ($h/\lambda\ll1$), where
the wake regenerates every step. An unconditional floor supplies
dissipation there, rolling off where the absorption length is resolved
so that its steady-state cost stays bounded:
\begin{equation}
  \alpha_{{\rm floor},i} = s_i\,
    \frac{\alpha_{\rm floor}}{1 + \big(h_i\,\kappa\rho_i/\varepsilon_\lambda\big)^4} ,
  \label{eq:fuv-diss-floor}
\end{equation}
with $\alpha_{\rm floor}$ (default $0.5$) its ceiling,
$\varepsilon_\lambda$ (default $0.5$) the value of $h/\lambda$ at which
the roll-off sets in, and $s_i\in[0,1]$ the relaxation-residual gate
below. At the default $\alpha_{\rm floor}$ and $\varepsilon_\lambda$,
the quartic tail leaves a coefficient of $3.1\times10^{-6}$ at
$h/\lambda\sim10$. $\alpha_{\rm floor} = 0$ recovers the trigger-only
scheme exactly.

\textbf{The flux-relaxation-residual gate.} Even with the roll-off, the
floor costs accuracy on a resolved, settled profile in the band
$0.25\lesssim h/\lambda\lesssim0.85$, where the physical absorption is
too weak to damp it and the front it protects has long passed. A settled
particle is recognisable: its flux sits at the fixed point of
Eq.~\ref{eq:fuv-fixed-point}. With $w_i = \kappa\rho_i + H/c$ -- the
\emph{true} speed of light, not $c_{{\rm hyp},i}$: $w_i$ is
$a_i/(c_{{\rm hyp},i}\Delta t_i)$ by construction, so $c_{{\rm hyp},i}$
cancels out of $w_i$ itself, unlike $a_i$ -- the residual
\begin{equation}
  R_i = \frac{\big|\,w_i\,\mathbf F_i + c_{{\rm hyp},i}\,\mathbf G_i\,\big|}
             {w_i\,|\mathbf F_i| + c_{{\rm hyp},i}\,|\mathbf G_i|}
  \in[0,1] , \qquad
  s_i = \min\!\left(1,\ \Big(\frac{R_i}{\varepsilon_R}\Big)^2\right) ,
  \label{eq:fuv-diss-gate}
\end{equation}
is zero exactly at the fixed point ($\mathbf F_i = -\mathcal C_i\mathbf
G_i$, the form above being that relation multiplied through by $w_i$ so
that no intermediate overflows at small $\kappa$), and close to one on
a fresh front, wherever only one of $\mathbf F_i$, $\mathbf G_i$ is
nonzero, or where $\tau\gg\Delta t$ so that the flux has not relaxed
yet. $\mathbf F_i$ here is the flux that produced $u^n$, i.e.\ the
previous step's limited flux. $\varepsilon_R$ (default
$0.40$) is the residual below which the floor is reduced; the gate can
only lower the floor, never raise it. Three rules cover the degenerate
inputs: a quiescent particle with $\mathbf F_i = \mathbf G_i = \mathbf 0$
has a zero denominator and gets $s_i = 0$ (it is trivially at the fixed
point and has nothing to damp); $w_i = 0$ (no absorption and no
expansion, hence no relaxation timescale) or $c_{\rm hyp} = 0$ gives
$s_i = 1$; $\varepsilon_R = 0$ disables the gate ($s_i = 1$
everywhere). A particle whose flux is pinned by the limiter ($f\approx1$)
never reaches the unlimited fixed point, so the gate keeps part of the
floor there: a conservative false positive (Section~\ref{sec:fuv-limits}).

\textbf{Application.} The exchange of Eq.~\ref{eq:fuv-diss-formula} is
accumulated in the force loop, with this step's coefficients, on $u^n$,
and decays with $u^n$ inside the energy update
(Eq.~\ref{eq:fuv-u-update}). The trigger reads $u^n$ against the kernel
mean of the neighbours' $u^n$, so an undershoot created by one step's
update is corrected by the next step's; cooling sees it for one step,
through its non-negative clamp. Folding the exchange in with the decay
factor $e^{-a_i}$, rather than adding it undamped, keeps the per-step
amplification matrix with the same eigenvalues as a correction applied
to an already relaxed field, which has the larger stability margin
(\texttt{verify\_isrf\_dissipation.py}, Part I); the discrete steady
state is then
\begin{equation}
  \big(1 - e^{-a_i}\big)\,u_i = \Delta t_i\,\varphi(a_i)
    \left[\frac{c_{{\rm hyp},i}}{c}\,\dot s_i
      - \Big(\tfrac{1}{\rho}\nabla\cdot(\rho\mathbf F)\Big)_i
      + e^{-a_i}\,\dot u_i\big|_{\rm diss}\right] .
  \label{eq:fuv-diss-apply}
\end{equation}
The transient amplification, $\max_{n\le200}\|G^n\|_2$ along the
closure, is not the same as for that placement: at $C_{\rm hyp} = 0.5$
it equals the reference ($1.35$) for $\alpha\lesssim0.323$, reaches
$1.54$ at the shipped $\alpha_{\max} = 0.5$, and $1.68$ at the guard
limit $\alpha = 0.5887$. At the guard boundary itself the ratio to the
reference grows with $C_{\rm hyp}$: $1.14\times$ at $C_{\rm hyp} =
0.25$, $1.25\times$ at $C_{\rm hyp} = 0.5$, and $1.35\times$ at
$C_{\rm hyp} = 1.0$
(\texttt{verify\_isrf\_dissipation\_flux\_first\_parameters.py}, Part E).
It is bounded within the stability bound of
Eq.~\ref{eq:fuv-stability-bound} (the spectral radius stays at or below
one there) and vanishes with the dissipation off. The dissipation conserves $\sum_i m_i u_i$ exactly
for pairs of active particles in the same time bin (Section~\ref{sec:fuv-operators});
it does not conserve the first spatial moment: behind a source moving through the
gas the trigger fires preferentially on the depleted trailing side, and
the field's centroid lags the source by a small additional amount,
measured at $0.2$--$1.8\%$ of $\lambda$ at the shipped ceiling. Because
the floor can act where the trigger never fires, the stability bound of
Eq.~\ref{eq:fuv-stability-bound} is checked on the larger of the two
ceilings.

\subsubsection{Source turn-off}
\label{sec:fuv-turnoff}

Once a star is no longer processed as an active feedback source it makes
no further deposits in either band (Section~\ref{sec:fuv-sources}).
With propagation on, each illuminated particle's reservoir drains to
zero by its horizon, at most one star step after the last visit, and
the field then obeys Eqs.~\ref{eq:fuv-u-eq}--\ref{eq:fuv-F-eq} with no
source: it spreads at $c_{\rm hyp}$ and decays with the $e$-folding time
$\tau = 1/\!\left[c_{\rm hyp}\big(\kappa\rho + H/c\big)\right]$, the two
decay rates (absorption and the $c_{\rm hyp}/c$-dilated expansion term)
adding directly (not their two $e$-folding times), so $\tau$ is shorter
than $1/(c_{\rm hyp}\kappa\rho)$ alone in a cosmological run -- but, with
$H$ correctly dilated, only by an amount suppressed by $c_{\rm hyp}/c$:
for any $\kappa\rho$ not itself many orders of magnitude below $H/c$,
the Hubble term's contribution to $\tau$ is negligible next to
absorption (Section~\ref{sec:fuv-lagrangian}).
The one-step drop of the source at a particle is the sharpest transient
the scheme sees, and the trigger may fire on the innermost particles for
a step or two before the field settles.

\subsubsection{Propagation off: the instantaneous field}
\label{sec:fuv-instantaneous}

With propagation off (the default), none of
Sections~\ref{sec:fuv-chyp}--\ref{sec:fuv-turnoff} runs. A star's
visit resets $u_{b,j}$ to zero on the first visit of the step and every
visit that step adds $\Delta E_{b,j}/m_j$ from Eq.~\ref{eq:fuv-inject};
the field is the energy deposited during the star's own step, confined
to the kernels of the stars active that step, and is zeroed when a
particle's illumination window lapses. Extinction still applies. This
mode is a diagnostic baseline that isolates injection, extinction and
the Grackle coupling. It is not an approximation to the propagated
field: it reaches no particle outside a kernel, and its magnitude is
proportional to the star's time step $\Delta t_\star$ rather than to the
absorption time that sets a physical energy density.

\subsubsection{Cosmological runs}
\label{sec:fuv-cosmo}

The fields, $c_{\rm hyp}$ and $\kappa\rho$ are physical; the positions,
separations, smoothing lengths and densities the loops read are comoving.
A comoving length is the physical one divided by $a$, a comoving density
the physical one multiplied by $a^3$, and a kernel derivative in three
dimensions carries $h^{-4}$, hence $a^{+4}$ on comoving input. Counting
term by term, each of the three pairwise operators returns exactly $a^{1}$
times the physical value it estimates:
\begin{itemize}
  \item \textbf{Flux divergence} (Eq.~\ref{eq:fuv-div-op}). The separation
    vector in $\mathbf F\cdot\nabla W$ gives $a^{-1}$ and the inverse
    separation inside the kernel gradient gives $a^{+1}$, which cancel;
    $1/\rho$ gives $a^{-3}$; the kernel derivative gives $a^{+4}$. Net
    $a^{1}$.
  \item \textbf{Pressure-tensor divergence} (Eq.~\ref{eq:fuv-grad-op}).
    The volumetric field $U = \rho u$ gives $a^{+3}$, the $m_j/\rho_i^2$
    normalisation $a^{-6}$, the kernel derivative $a^{+4}$. Net $a^{1}$.
  \item \textbf{Artificial dissipation} (Eq.~\ref{eq:fuv-diss-formula}).
    The jump $U_i - U_j$ gives $a^{+3}$, the $1/(\rho_i\rho_j)$ prefactor
    $a^{-6}$, the symmetrised kernel derivative $a^{+4}$. Net $a^{1}$.
\end{itemize}
Each is therefore closed by one division by $a$, applied where the
comoving estimate is written into its physical accumulator, and nothing
else inside the operators needs a conversion. The exception is the
neighbour reference field $\langle|U|\rangle_i$ of the dissipation
trigger, which takes no conversion at all: it is read only as a ratio
against $U_i$, which carries the same weight, so converting it would
introduce a bias rather than remove one.

The smoothing length in Eq.~\ref{eq:fuv-chyp} and in the dissipation
formulas is the physical one, $a\,h$; $\Delta t$ is proper time; the
expansion enters only through the $H$ terms derived in
Section~\ref{sec:fuv-lagrangian}, folded into the relaxation depth of
Eq.~\ref{eq:fuv-phi} so that the homogeneous decay stays exact at any
$H\Delta t$. The injection column of Eq.~\ref{eq:fuv-extinction} is
converted to physical units before the exponential.

\subsubsection{Runtime parameters}
\label{sec:fuv-parameters}

All keys live in the \code{GEARFeedback} block. The photoelectric
switch is the master switch; the rest are read only when it is on and
are meaningful only with propagation on.

\begin{center}
\small
\begin{tabular}{|p{4.2cm}|p{1.3cm}|p{1.3cm}|p{1.9cm}|p{4.6cm}|}
\hline
\textbf{Key} & \textbf{Symbol} & \textbf{Default} & \textbf{Range} &
  \textbf{Meaning} \\
\hline
\code{with\_photoelectric\_heating} & & 0 & 0 or 1 &
  Master switch: compute $L_{\rm FUV}$, $L_{\rm LW}$, inject, and couple
  $G_0$ and $k_{\rm diss}$ into Grackle. \\
\hline
\code{photoelectric\_heating\_efficiency} (\code{GrackleCooling}) & &
  \code{constant} & \code{constant}, \code{wolfire1995},
  \code{density\_dependent} &
  Grackle photoelectric efficiency (Eq.~\ref{eq:fuv-pe-heating} and
  Eq.~\ref{eq:fuv-pe-efficiency-density}). \\
\hline
\code{ISRF\_propagation} & & 0 & 0 or 1 &
  0: instantaneous field (Section~\ref{sec:fuv-instantaneous}); 1:
  two-moment propagation with reservoir, transport and dissipation. \\
\hline
\code{ISRF\_c\_hyp\_margin} & $C_{\rm hyp}$ & 0.5 &
  $(0, 1.69]$, and Eq.~\ref{eq:fuv-stability-bound} &
  Courant number of the propagation, $c_{\rm hyp} = \min(C_{\rm hyp}
  h/\Delta t, c)$. \\
\hline
\code{ISRF\_dissipation\_alpha\_max} & $\alpha_{\max}$ & 0.5 &
  $\ge0$, Eq.~\ref{eq:fuv-stability-bound} &
  Ceiling of the negativity-triggered coefficient. 0 disables the
  trigger. \\
\hline
\code{ISRF\_dissipation\_negativity\_threshold} & $\varepsilon_1$ & 0.01 &
  $(0, 1]$ &
  Relative undershoot at which the trigger saturates. \\
\hline
\code{ISRF\_dissipation\_alpha\_floor} & $\alpha_{\rm floor}$ & 0.5 &
  $\ge0$, Eq.~\ref{eq:fuv-stability-bound} &
  Ceiling of the unconditional floor. 0 disables the floor. \\
\hline
\code{ISRF\_dissipation\_floor\_h\_over\_lambda} & $\varepsilon_\lambda$ &
  0.5 & $>0$ &
  $h/\lambda$ at which the floor's quartic roll-off sets in. \\
\hline
\code{ISRF\_dissipation\_floor\_relaxation\_residual} & $\varepsilon_R$ &
  0.40 & $[0, 1]$ &
  Residual below which the floor is reduced by $(R/\varepsilon_R)^2$.
  0 disables the gate. \\
\hline
\code{ISRF\_c\_hyp\_pin\_for\_debugging} & & 0 & $\ge0$ &
  Test only: pin every particle's $c_{\rm hyp}$ to this value,
  bypassing the clamp at $c$. Never in production. \\
\hline
\code{ISRF\_dissipation\_alpha\_pin\_for\_debugging} & & 0 & $\ge0$ &
  Test only: hold every particle's coefficient at this value, bypassing
  trigger and floor. Never in production. \\
\hline
\end{tabular}
\end{center}

\subsection{Validation summary}
\label{sec:fuv-validation}

The scheme is exercised by a battery of single-star and seeded-pulse
tests in periodic glass boxes, each isolating one regime, plus unit
tests of the operators. The regimes, the pass criteria and the outcome
at the shipped defaults are summarised here; the full record, with the
fixtures, the numbers and their history, is in the GEAR radiation
logbook.

\begin{itemize}
  \item \textbf{Injection normalisation.} Propagation off, $Z=0$ so
    that $\zeta = 1$: the energy summed over the star's kernel must
    equal $\Delta t_\star L_b$ to floating-point precision, per band.
  \item \textbf{Time-bin seam conservation.} A fixture straddling a
    time-bin boundary, dissipation off, $Z=0$, checked per band against
    a uniform-glass control carrying the same dose in flight:
    $E(t) = (c/c_{\rm hyp})\sum_i m_i u_i/L - t$. Criterion:
    $|E_{\rm seam}(t) - E_{\rm uniform}(t)|/t \le 10^{-5}$ at every
    snapshot $t>0$. Outcome: $9.3\times10^{-8}$ (FUV), $6.6\times10^{-8}$
    (LW), down from $5.0\times10^{-2}$ before the pair with
    $H_i \le r_{ij} < H_j$ was dispatched from the force loop
    (Section~\ref{sec:fuv-operators}).
  \item \textbf{Screened steady state}, $h/\lambda$ from $0.6$ to $20$
    (production particle masses $0.1$--$760\,M_\odot$ at solar $Z$,
    $n_H = 10^3\,{\rm cm^{-3}}$). Criterion: the settled radial profile
    $u(r)$, binned and compared bin by bin against the scheme's own
    discrete fixed point (computed on the run's particle positions, in
    both bands), median relative error below $15\%$ over every radial
    bin whose predicted value clears a float32-noise floor
    ($10^3\times$ the smallest normal float32); a bin below that floor,
    or a band with fewer than three such bins, is reported unmeasurable
    rather than counted as a pass or a fail; the amplitude identity
    Eq.~\ref{eq:fuv-amplitude-identity} reported alongside. Outcome:
    pass at every corner, with amplitude agreement to $0.2\%$ or
    better; the LW profile at $h/\lambda = 20$ is unmeasurable because
    the field is zero on all but a few dozen particles. The
    $95\,M_\odot$ (LW) and $760\,M_\odot$ (FUV) corners exercise the
    float32 underflow of Section~\ref{sec:fuv-moments} directly: median
    relative error $0.5014$ and $0.5551$ before evaluating the flux norm
    in double, $0.0074$ and $0.0104$ after, on the same check.
  \item \textbf{Causal reach.} The field's edge at $10^{-3}$ of the peak
    must stay inside $r = c_{\rm hyp}(t - t_0)$ at every snapshot, in a
    uniform box and with one neighbour pinned to a $c_{\rm hyp}$ up to
    $46\times$ its neighbours'. Outcome: pass in every case.
  \item \textbf{Optically thin corner}, $h/\lambda\sim6\times10^{-3}$,
    and the \textbf{transient ladder}, $h/\lambda$ from $0.0028$ to
    $0.42$ with the front still inside the box. Criterion: the worst
    radial bin's negative-mass fraction at or below $0.05\%$ at the end
    of the run. Outcome: pass at every rung; without the floor the
    fraction rises to $\sim30\%$ at the thinnest rungs, and without any
    dissipation to $\sim60$--$80\%$.
  \item \textbf{Transient fronts in the band} $h/\lambda = 0.15$,
    $0.25$, $0.35$, both bands. Same criterion, plus agreement within
    $0.02$ percentage points with the ungated floor. Outcome: $0.0000\%$
    at all three rungs.
  \item \textbf{Settled profiles in the band} $h/\lambda = 0.25$ to
    $0.85$. Criterion: median relative error of $u(r)$ against the
    discrete fixed point at or below $0.15$. Outcome: $0.12$, $0.12$,
    $0.12$, $0.085$, $0.030$ at $h/\lambda = 0.25$, $0.35$, $0.45$,
    $0.61$, $0.85$; the ungated floor fails this criterion at the first
    four points (up to $0.42$). Transient negativity peaks of at most
    $0.12\%$ during front passage settle to exactly zero.
  \item \textbf{Production point}, $95\,M_\odot$ particles, $n_H =
    1\,{\rm cm^{-3}}$, $Z = 10^{-2}Z_\odot$ ($h/\lambda = 6\times10^{-4}$).
    Outcome: worst-bin negative-mass fraction $0.007\%$ late in the run,
    $0.027\%$ at the single worst snapshot, against $17\%$ with the
    trigger alone.
  \item \textbf{Source turn-off.} A $70\,M_\odot$ star at solar $Z$
    dying at $3.97$\,Myr. Criterion: the post-death $e$-folding time of
    the box-total field equals $1/(c_{\rm hyp}\kappa\rho)$, and no
    post-death snapshot carries a negative field. Outcome: ratio $1.000$
    (FUV) and $0.999$ (LW), $0$ negative snapshots of $67$ in either
    band.
  \item \textbf{Unit tests.} Conservation of the dissipation exchange
    across a smoothing-length ratio, scale-factor independence of the
    three spatial operators, and the sign and magnitude of the expansion
    term. All pass.
  \item \textbf{End-to-end heating.} A heating-on versus heating-off
    pair through Grackle checks the sign and magnitude of the thermal
    response; no quantitative accuracy target is defined for the $G_0$
    path beyond that.
\end{itemize}
Restart from a mid-run dump reproduces an uninterrupted control run;
the per-particle state is covered by the existing whole-struct restart
path.

\subsection{Limitations and open items}
\label{sec:fuv-limits}

\begin{itemize}
  \item \textbf{Limiter-bound residual on settled profiles.} Particles
    whose flux is pinned by the limiter ($f\approx1$, most of the shell
    $0.5\lesssim r/\lambda\lesssim1$ around a single source) never reach
    the unlimited fixed point of Eq.~\ref{eq:fuv-fixed-point}, so the
    gate retains $\sim40\%$ of the floor there in a fully settled
    state. The shape error in the $0.25$--$0.85$ band passes its
    criterion but stays $\sim6\times$ above what a uniform, tiny floor
    would give at the same point. Measuring $R$ against the limited
    recurrence's own fixed point is the candidate improvement.
  \item \textbf{Untested branches.} The gate and the $H$ terms have
    been unit-tested for scale-factor independence but not exercised in
    a cosmological run with $a\ne1$ and $H\ne0$. The many-source,
    near-isotropic branch (Eq.~\ref{eq:fuv-steady-diffusion}), which is
    the production use, has not been validated directly; every
    propagation test above uses one source or one seeded pulse.
    Production resolution has been tested at one point only.
  \item \textbf{Symmetric divergence dispatch: a positivity cost.}
    Reaching every pair with $H_i \le r_{ij} < H_j$ from the force loop
    (Section~\ref{sec:fuv-operators}) fixed the time-bin seam's
    energy-conservation error above, but is not free. In a strong
    hot/cold density-contrast fixture (a $10^6$\,K neighbour against a
    cold background) the early-transient negative-$u$ population grows
    to $323$ particles (FUV) and $456$ (LW) at snapshots $3$ to $8$,
    against at most $1$ before the dispatch moved; a smaller, standing
    negative population also appears at time-bin boundaries, where
    particles that do not share a time bin or the same $c_{\rm
    hyp}\kappa\rho$ leave a residual by construction
    (Section~\ref{sec:fuv-operators}). Both are clamped to zero before
    Grackle sees them (Section~\ref{sec:fuv-coupling}); improving
    positivity at time-bin boundaries and in strongly contrasted gas,
    without changing transport in smooth fields, is open. This is a
    different mechanism from the $c_{\rm hyp}$-heterogeneity pathway of
    Section~\ref{sec:fuv-chyp-heterogeneity} (a dispatch-coverage gap
    versus a receiver-side Courant violation), and the two should not be
    conflated: fixing one does not address the other.
  \item \textbf{Dissipation coefficients.} The amplification matrix of
    the flux-first update (Eq.~\ref{eq:fuv-u-update}) has the same
    characteristic polynomial as a dissipation applied to the transported
    field, and the thin limit $e^{-a}\to1$ binds, so
    Eq.~\ref{eq:fuv-stability-bound} is the stability condition of the
    shipped order: at $C_{\rm hyp} = 0.5$ it admits $\alpha\le0.589$,
    $1.18$ times $\alpha_{\max} = 0.5$. The bound uses $2I_W = 6.2$; the
    lattice maximum of the dissipation symbol is $3.31$, so on a uniform
    lattice the bound is about twice conservative in $\alpha$. A direct
    grid scan of the discrete recurrence's instability onset in
    $C_{\rm hyp}$, at fixed $\alpha \in \{0.25, 0.5, 0.75\}$, places it
    at $1.97$--$2.58$ times the analytic root of
    Eq.~\ref{eq:fuv-stability-bound}, consistent with that same margin
    rather than with any unmodelled effect. The transient
    amplification of the shipped order grows with $\alpha$ above
    $\alpha\approx0.32$: $1.54$ at $0.5$ and $1.68$ at the bound, against
    $1.35$ with the dissipation off
    (\texttt{verify\_isrf\_dissipation\_flux\_first\_parameters.py}). The
    negativity threshold $\varepsilon_1 = 0.01$ only sets where the
    trigger saturates and does not enter the stability analysis; it was
    measured under the isotropic closure and not re-measured with the
    limiter and floor. $\varepsilon_\lambda = 0.5$ is the least
    constrained of the dissipation parameters.
  \item \textbf{Moment-closure physics.} Shadows behind dense clumps are
    softer than reality, and two beams crossing at a particle are merged
    into one direction; both are intrinsic to a two-moment closure. A
    source moving relative to its gas trails a field lagging by
    $v_{\rm rel}\lambda/c_{\rm hyp}$, $c/c_{\rm hyp}$ times the physical
    lag.
  \item \textbf{Instantaneous mode.} The propagation-off field scales
    with the star's time step and reaches only kernel neighbours; it is
    a baseline, not a physical field.
  \item \textbf{Grackle path.} The photoelectric efficiency is constant,
    follows the grain charging parameter with Grackle's electron density
    \citep[their Eq.~2]{Wolfire1995}, or is a function of density alone
    (Eq.~\ref{eq:fuv-pe-efficiency-density}). The density-only form
    assumes gas on the solar-neighbourhood thermal equilibrium curve and
    overestimates the efficiency at low metallicity \citep{Smith2026}.
    The default is not yet chosen on evidence. The H$_2$ cross-section is an effective
    band average with an implicit spectral shape; no quantitative
    accuracy target exists yet for the $G_0$ path.
  \item \textbf{Multi-source excess noise at high resolution, open.} A
    lattice of many co-located sources (\texttt{ISRFMultiSourceSuperposition}),
    checked outside every star's own kernel to probe the many-source
    branch flagged untested above, agrees with the continuum
    free-streaming prediction at a coarser gas resolution and does not at
    a finer one: the second moment of $u$ runs $0.349$/$0.360$
    (root-mean-square of $u/\langle u\rangle-1$, FUV/LW) against a
    continuum prediction of $0.028$/$0.047$, about $12\times$ high in FUV
    and $8\times$ in LW, and keeps changing from snapshot to snapshot
    rather than settling, while the first moment $\langle u\rangle$
    agrees with the prior implementation to four to six significant
    figures. At the finer resolution the flux is pinned by the limiter
    for essentially all of that outside-kernel gas ($99.75\%$ FUV,
    $98.84\%$ LW), so the flux-relaxation-residual gate of
    Eq.~\ref{eq:fuv-diss-gate} saturates at $s=1$ there, for the reason
    already given in Section~\ref{sec:fuv-dissipation}: a limiter-pinned
    flux never reaches the unlimited fixed point. That saturation is neither
    necessary nor sufficient for the excess noise: removing the floor,
    and removing the floor and the trigger together, both still fail the
    fixture and make every other diagnostic worse (the second moment,
    the step-to-step variation of $u$, and positivity all degrade
    further), while forcing the same saturated coefficient onto the
    coarser resolution leaves it passing and improves every diagnostic
    there. The excess noise's own cause remains under investigation.
\end{itemize}
